\documentclass{article}

% if you need to pass options to natbib, use, e.g.:
%     \PassOptionsToPackage{numbers, compress}{natbib}
% before loading neurips_2026

\usepackage{neurips_2026}
% Switch to [eandd] for Datasets \& Benchmarks track or [preprint] for arXiv.

\usepackage[utf8]{inputenc}
\usepackage[T1]{fontenc}
\usepackage{hyperref}
\usepackage{url}
\usepackage{booktabs}
\usepackage{amsfonts}
\usepackage{amsmath}
\usepackage{nicefrac}
\usepackage{microtype}
\usepackage{xcolor}
\usepackage{multirow}
\usepackage{listings}
\usepackage{graphicx}
\usepackage{fvextra}
\usepackage{tabularx}


% TODO / placeholder macros
\newcommand{\todo}[1]{\textcolor{red}{\textbf{TODO: #1}}}
\newcommand{\lianhui}[1]{\textcolor{red}{[#1 --LQ]}}
\newcommand{\tbd}{\textcolor{red}{\textbf{--}}}

% Verbatim include with linewrapping for long XML/transcript files.
\newcommand{\promptinput}[2][]{%
  \VerbatimInput[
    fontsize=\footnotesize,
    breaklines=true,
    breakanywhere=true,
    breaksymbolright=\small$\hookrightarrow$,
    obeytabs=true,
    tabsize=2,
    frame=leftline,
    framerule=0.6pt,
    rulecolor=\color{gray},
    #1
  ]{#2}%
}
\usepackage{xspace}

% Earlier title variants (kept for reference):
% \title{Coding Agent Adaptation for Advanced Scientific Tooling: An Application Study in Automating Geophysics Simulations}
% \title{Simulator-Interface Grounding Adapters for Scientific Simulation Setup: A Geophysics Case Study}
\title{Adapting self-evolving Off-the-Shelf Coding Agents for Scientific Simulator Setup: A GEOS Case Study}


\author{%
  Anonymous Authors\\
}


\begin{document}


\maketitle


\begin{abstract}
% Advanced simulators and other powerful scientific software ship with correspondingly expressive configuration languages, and learning to operate them costs domain scientists hours to days per study. This \emph{advanced tool-operating} bottleneck is a concrete, near-term target for agentic AI, and a precondition for trusting agents with the open-ended scientific reasoning that has drawn most current attention. We study how far this bottleneck can be reduced by adapting an off-the-shelf coding agent (Claude Code) into a simulator-operating assistant, taking the open-source GEOS multiphysics simulator (used in CO$_2$-storage and induced-seismicity research) as our case study. Rather than build a bespoke agent loop, we wrap the harness with four grounding components, each addressing a known agent failure mode: retrieval over simulator artifacts (R), a schema-validating stop-hook (S), an agent-callable XML validator (X), and a procedural-memory cheatsheet (M). We call this lightweight adaptive approach a \textbf{Simulator-Interface Grounding Adapter (SIGA)}, and present detailed ablations to dissect which components carry the empirical weight. On a held-out set of compound multi-physics tasks, our adapters lift mean TreeSim by $+0.07$ and reduce across-seed standard deviation by roughly an order of magnitude over the bare agent, with attribute-level semantic errors remaining as the residual frontier. A human baseline of geoscience-domain experts new to GEOS places the SIGA agent inside the time envelope of an experienced human author and at parity with extended-budget human deck quality on a representative task. We also explore agent autonomy by varying the level of detail of the task specification and equipping the agent with a tool to consult a human supervisor when ambiguities arise. Finally, on top of this constrained design space we show that an open-ended self-evolution pipeline matches the best hand-designed cell, and a pilot OpenFOAM transfer reproduces the dominant mechanism, suggesting the recipe is not GEOS-specific. 

Advanced scientific simulators expose specialized input languages that turn simulation goals into executable configurations, but learning these interfaces can cost domain scientists hours to days. We study simulator setup as a problem of agent-tool interface grounding: what minimal simulator-specific contract is needed for an off-the-shelf coding agent to operate real scientific software? Our intuition is that coding agents already know how to navigate files, edit code, run commands, and repair outputs, but they lack the simulator’s executable contract: its vocabulary, structural constraints, validation rules, and termination conditions. We introduce SIGA, a Simulator-Interface Grounding Adapter that supplies this contract through retrieval, procedural memory, in-trajectory validation, and validation-enforced termination. On GEOS deck authoring, SIGA converts a representative Buckley–Leverett setup from an hours-scale expert workflow into a roughly five-minute agent task, reaching deck-level TreeSim above 0.90 at parity with extended-budget expert quality---a ${\sim}36\times$ speedup over a domain expert new to GEOS. On a harder held-out set, grounding raises mean TreeSim by roughly $7$ points (${\sim}10\%$ relative, $0.720 \to 0.789$) and cuts across-seed variance by about an order of magnitude, at no added wall-clock cost over the bare agent; the best held-out adapter is found automatically by self-evolution. OpenFOAM and LAMMPS transfers show that the useful grounding mechanism shifts by interface: validation improves completeness-heavy simulators, while memory and retrieval improve domain-correct script authoring.

\end{abstract}


\section{Introduction}
\label{sec:intro}
\begin{figure}[t]
\centering
\includegraphics[width=\textwidth]{assets/geos_fig1_v3(1).png}
\caption{Illustrative example of advanced tooling usage bottleneck for the geophysics domain. Here, the GEOS simulator's extensive documentation helps as a translation guide for its elaborate configuration that is a custom XML (domain specific language/DSL) to produce results such as simulating carbon sequestration in deep saline formations (top right) or reservoir flow in heterogeneous hydrocarbon (bottom right).}
% \caption{\textbf{Comparison of manual and Simulator-Interface Grounding Adapter (SIGA) workflows.} While the manual workflow \textbf{(a)} requires hours of iteratively navigating documentation and debugging, the SIGA workflow \textbf{(b)} reduces setup to minutes by supplying an LLM agent with structured inputs and specialized tools. Both approaches \textbf{(c)} ultimately generate a valid XML deck to execute complex, coupled multi-physics simulations in the GEOSX engine.}
\label{fig:1}
\end{figure}

% Numerical simulation is among the most widely used computational workflows in modern science, supporting research across physics, chemistry, biology, materials science, climate, biomedicine, and the geosciences. Yet between a researcher's natural-language intent and a runnable simulation lies a substantial expert bottleneck: operating any advanced piece of scientific software amounts to learning its domain-specific language (DSL). A simple tool can be operated through a handful of function calls, while a powerful one demands a correspondingly expressive DSL whose vocabulary corresponds to objects, methods, and parameters in the underlying codebase. For modern multiphysics simulators this DSL takes the form of input decks, namelists, or configuration scripts whose tag names refer to specific simulator classes, whose attribute values must satisfy schema and physical constraints, and whose cross-references must remain mutually consistent across many sections. Configuring such software routinely costs trained domain scientists hours to days per study. Reducing this DSL-operating bottleneck would have direct impact across scientific fields, and is a necessary precondition for entrusting agents with more open-ended scientific responsibilities further down the line.
% Recent work has explored LLM agents for scientific software, with system-specific agents proposed for chemistry~\citep{bran2024chemcrow,boiko2023coscientist}, computational fluid dynamics~\citep{chen2024metaopenfoam,yue2025foamagent,pandey2025openfoamgpt}, molecular dynamics~\citep{shi2026mdagent2,zhao2026polyjarvis,guilbert2025dynamate}, finite-element and multiphysics simulation~\citep{zhan2025mooseagent,mcpsim2026,ni2024mechagents}, and reservoir simulation~\citep{moyner2026jutulgpt}. Two complementary gaps motivate our study. \emph{Application-wise}, no agent has been designed for GEOS in particular, despite its central role in subsurface science: existing geoscience language models focus on knowledge access and question answering~\citep{deng2023k2,lin2024geogalactica}, and the closest geoscience-agent systems target narrower simulator families~\citep{bekele2025geosim,li2025seismologyagent,moyner2026jutulgpt}. \emph{Method-wise}, prior scientific-agent work largely re-implements the orchestration layer from scratch (custom agent loops, tool-calling logic, retry and termination handling), rather than building on the substantial engineering already invested in frontier coding agents such as Claude Code or open-source efforts such as OpenCode and OpenHands~\citep{wang2024openhands}. As these general coding agents become competent on their own, a different starting point becomes available: take a SOTA coding harness as given, and invest engineering only in the simulator-specific grounding layer above it. The empirical question this raises is what that grounding layer needs to contain to make a generic coding agent reliable on a real scientific-software target.


Scientific simulation is one of the central computational workflows in modern science. Researchers use simulators to study physical, chemical, biological, and engineering processes that would be expensive, slow, dangerous, or impossible to observe directly in the laboratory. In geophysics, simulators such as the GEOS \footnote{\url{https://github.com/GEOS-DEV/GEOS}}~\citep{geos2024} model CO$_2$ sequestration, reservoir flow, hydraulic fracture, wellbore behavior, and induced seismicity; in molecular science, molecular-dynamics engines such as LAMMPS \citep{holbrook2026lammps} encode atomistic systems, force fields, and time-integration protocols. Across these domains, the scientific goal may be expressed in natural language, but the executable simulation must be written in a simulator-specific input language.

This translation from scientific intent to runnable simulator configuration is a persistent bottleneck in simulation-driven research. Advanced simulators are not simple tools with a few stable arguments. They expose large domain-specific language (DSL) whose tokens name solver classes, material models, mesh regions, boundary conditions, event schedules, numerical schemes, and output requests. A valid input deck must satisfy syntax constraints, schema constraints, and physical or domain-conventional constraints, while maintaining consistent names and references across sections and often across files. Scientists new to a simulator can therefore spend hours or days reading documentation, searching examples, testing numerical settings, and debugging invalid configurations before obtaining a runnable deck.

This setup stage is a natural target for agentic AI because it is both valuable and constrained. It is valuable because reducing setup time directly accelerates scientific iteration. It is constrained because the agent is not being asked to invent a new hypothesis or interpret final physical results; it is being asked to operate an existing scientific tool according to a given specification. This makes simulator setup a useful near-term testbed for scientific agents. If current agents cannot reliably translate a specification into a simulator DSL, they are unlikely to be dependable in broader workflows that require experimental design, hypothesis revision, and domain reasoning.

Recent scientific-agent systems have explored chemistry~\citep{bran2024chemcrow,boiko2023coscientist}, molecular dynamics~\citep{shi2026mdagent2,zhao2026polyjarvis,guilbert2025dynamate}, computational fluid dynamics~\citep{chen2024metaopenfoam,yue2025foamagent,pandey2025openfoamgpt}, finite-element and multiphysics modeling~\citep{zhan2025mooseagent,mcpsim2026,ni2024mechagents}, reservoir simulation~\citep{moyner2026jutulgpt}, and other simulator-centered workflows. Two complementary gaps motivate our study. \emph{Application-wise}, no agent has been designed for GEOS in particular, despite its central role in subsurface science: existing geoscience language models focus on knowledge access and question answering~\citep{deng2023k2,lin2024geogalactica}, and the closest geoscience-agent systems target narrower simulator families~\citep{bekele2025geosim,li2025seismologyagent,moyner2026jutulgpt}. \emph{Method-wise}, many of these systems build simulator-specific agent loops from scratch, including custom planning, tool calling, retry logic, and termination handling, rather than building on the substantial engineering already invested in frontier coding agents such as Claude Code or open-source efforts such as OpenHands~\citep{wang2024openhands}. We study a different design point: adapting an off-the-shelf coding agent. Modern coding agents already provide useful infrastructure for file navigation, iterative editing, shell use, and code repair. The remaining question is what simulator-specific grounding layer is needed to make such agents reliable simulator setup assistants.

We answer this question with the Simulator-Interface Grounding Adapter (SIGA), a thin wrapper around an off-the-shelf coding agent. The design is based on a simple intuition: modern coding agents already provide useful general machinery for navigating files, editing code, running commands, and repairing outputs, but they do not know the exact contract imposed by a scientific simulator. In simulator setup, failures therefore arise less from the lack of a generic agent loop than from missing interface grounding: the agent may not know the simulator’s vocabulary, may drift away from schema or structural constraints while editing, or may stop after producing a plausible-looking but incomplete deck. SIGA addresses this by grounding the agent at the points where these failures occur: before generation, during self-correction, and at termination.

Concretely, SIGA factorizes this grounding into four reusable components. Retrieval gives semantic access to simulator documentation, schemas, examples, and technical snippets; procedural memory keeps high-frequency simulator vocabulary and configuration patterns always visible; an agent-callable validator lets the agent check and repair outputs mid-trajectory; and a stop-hook makes validation a mandatory termination condition, preventing the agent from finishing with incomplete or structurally invalid deliverables. This decomposition is simulator-agnostic: the same slots instantiate as XML schema checks for GEOS, required-file and dictionary checks for OpenFOAM, and checks over commands, references, force-field declarations, and run-control blocks for LAMMPS. Because SIGA is a small adapter rather than a full agent loop, it also creates a compact optimization surface: we can ask whether prior trajectories can automatically improve the adapter’s primer, memory, and auxiliary skills.

Our experiments show that simulator-interface grounding compresses simulator setup from an hours-scale expert workflow into a roughly five-minute agent task while simultaneously improving deck quality and reliability. On a representative Buckley–Leverett GEOS deck-authoring problem, a geoscience-domain expert new to GEOS reached a complete, high-quality deck (deck-level TreeSim $0.931$) only after roughly three hours, whereas the SIGA-equipped agent matched that quality (deck-level TreeSim above $0.90$, where TreeSim is a 0–1 tree-similarity metric for deck quality) in about five minutes---a ${\sim}36\times$ reduction in wall-clock. The same grounding also helps relative to the bare agent: on this representative task it lifts deck-level TreeSim from $0.751$ (vanilla Claude Code) to above $0.90$, and on a harder held-out task set it raises mean TreeSim by about $7$ points ($0.720 \to 0.789$, a ${\sim}10\%$ relative gain) while shrinking across-seed standard deviation by roughly an order of magnitude ($0.081 \to 0.005$)---all at no additional wall-clock cost over the bare agent. The best held-out configuration is moreover discovered automatically: a self-evolved variant that rewrites its grounding components from prior trajectories attains this best performance. Supporting ablations and transfer studies explain how the useful mechanism shifts with the simulator interface: validation enforces structural completeness in GEOS and OpenFOAM, while memory and retrieval provide domain knowledge in LAMMPS, where the best adapter improves mean judge score from 4.56 to 7.78 on a 0–10 scale.

Taken together, these results suggest that near-term AI-for-science agents can deliver value by mastering the interfaces of existing scientific tools, rather than by replacing scientific reasoning. A small, self-improvable grounding layer around a general coding agent can compress a representative deck-authoring task from hours to minutes, while adapting its mechanism to the simulator interface: validation when completeness is the bottleneck, and memory or retrieval when domain correctness is the bottleneck.

% We address this question with the \textbf{Simulator-Interface Grounding Adapter (SIGA)}, a lightweight wrapper around an off-the-shelf coding agent (Claude Code) that turns it into a reliable simulator-operating assistant. We take the open-source GEOS multiphysics simulator\footnote{\url{https://github.com/GEOS-DEV/GEOS}}~\citep{geos2024} as our case study; GEOS is widely used in subsurface science for problems such as CO$_2$-storage fault-slip risk~\citep{he2026faultslip}, history-matching surrogates~\citep{han2024accelerated}, and microseismic-cloud mapping~\citep{glubokovskikh2023microseismic}, and is configured through XML decks spanning ten canonical sections that together form its operating DSL (Fig.~\ref{fig:1}; details deferred to \S\ref{sec:background}). SIGA wraps the harness with four grounding components, each instantiating an idea the broader AI literature has individually validated: semantic retrieval over simulator artifacts, as an additional search avenue beyond the agent's built-in filesystem tools; self-refinement baked into the harness's control flow as a schema-validating termination hook; the same validation tool also made available at the agent's discretion for mid-trajectory self-checking; and a procedural-memory cheatsheet distilled from prior trajectories and injected as always-on system-prompt context. Our main experiment is a detailed factorial ablation that disambiguates which of these components carries the empirical weight. We additionally study a self-evolved variant in which an offline pipeline rewrites the adapter contents against held-out training-set scores, testing whether automatic discovery in this design space is tractable.

% The empirical study identifies reliability, not raw quality, as the axis on which our adapters help. On in-distribution validation tasks the bare coding agent already performs near a quality ceiling and no component is reliably distinguishable from seed noise; on a harder held-out set the same components substantially reduce zero-score failures, lifting mean TreeSim by $+0.07$ and reducing across-seed standard deviation by roughly an order of magnitude. A bottleneck decomposition shows schema-aware adapters reliably fix block-level omissions while attribute-level semantic errors are untouched, identifying the residual frontier. A human baseline of geoscience-domain experts new to GEOS places the agent inside the time envelope of an experienced human author and at parity with extended-budget human deck quality on a representative task. A pilot cross-simulator transfer to OpenFOAM reproduces the dominant mechanism. Perhaps the most surprising finding is that even in this well-constrained, well-defined version of the task (which we view largely as a busy-work translation problem that does not demand substantial scientific reasoning, since the agent is given a detailed specification and only needs to express it in the simulator's DSL), SOTA coding agents on their own are not quite sufficient. This motivates both our study and continued investment in agentic systems for these workflows: even the more constrained bottleneck of operating scientific tools is unreliable with current systems, and improving reliability here provides immediate value to our domain partners while harder forms of scientific reasoning remain out of reach.



%% ============================================================
\section{Related work}
\label{sec:related}

\paragraph{LLM agents for scientific code and simulators.}
Scientific-coding agents are studied on curated benchmarks (ScienceAgentBench, DA-Code~\cite{chen2024scienceagentbench,huang2024dacode}) and built as general-purpose harnesses (OpenHands, SWE-agent~\cite{wang2024openhands,yang2024sweagent}), self-debugging loops~\cite{chen2024selfdebug}, or specialised search frameworks (SciNav~\cite{zhang2026scinav}). Domain-application systems include ChemCrow~\cite{bran2024chemcrow} and Coscientist~\cite{boiko2023coscientist} for chemistry, CellVoyager~\cite{alber2025cellvoyager} for single-cell analysis, and The AI Scientist~\cite{lu2024aiscientist,yamada2025ai} for end-to-end ML research. Closer to our setting are agents for scientific simulators: OpenFOAM-oriented computational fluid dynamics agents~\cite{chen2024metaopenfoam,pandey2025openfoamgpt,foamagent2025}, molecular-dynamics agents~\cite{kim2024mdagents,shi2026mdagent2,guilbert2025dynamate,zhao2026polyjarvis}, finite-element/mechanics agents~\cite{zhan2025mooseagent,ni2024mechagents}, and reservoir-simulation agents~\cite{moyner2026jutulgpt}. The shared design lesson is that general-purpose LLMs are insufficient for high-fidelity simulator use unless wrapped with domain grounding, structured interfaces, execution feedback, and iterative correction. Unlike most of these we build \emph{on top of} an existing engineered coding harness (Claude Code) rather than writing an agent loop from scratch. Foam-Agent 2.0~\cite{foamagent2025} is the most direct comparator for our OpenFOAM transfer study (\S\ref{subsec:openfoam-transfer}); we compare against its lint-only execution mode and treat that as a baseline.

\paragraph{Self-evolving agents, memory, and procedural guidance.}
Recent agent work treats capability as an iterative loop of planning, execution, feedback, and refinement that externalises experience into memory, search states, or revised plans~\citep{bran2024chemcrow,boiko2023coscientist,alber2025cellvoyager,lu2024aiscientist,zhang2026scinav,li2026agentflow,tang2026eigenagent,narayanan2024aviary}. A particularly active subfield, which treats the agent's own scaffolding as a learnable object (variously studied under the headings of meta-harness design, harness-as-code, agentic harness engineering, and skill optimization), has shown promising results on core AI benchmarks (coding, terminal navigation, math). Our self-evolved variant (\S\ref{subsec:self-evolved}) is fundamentally similar to these methods in that the agent reflects on its prior trajectories and rewrites its own plugin (the adapter); we contribute an initial study of how this paradigm fares on a task driven by domain knowledge rather than by general programming competence. Adjacent to this is a line on procedural memory and cheatsheets for LLMs (Buffer of Thoughts~\cite{yang2024bot}); we return to a specific finding about the procedural-memory interface in \S\ref{sec:discussion}.


%% ============================================================
\section{Background: GEOS as a domain-specific simulator language}
\label{sec:background}

A GEOS deck is one or more XML files that specify a multiphysics simulation across ten canonical sections, covering the mesh, geometry, execution schedule, physics modules, material models, computational regions, numerical methods, field specifications, functions, and outputs. Although the surface syntax is XML, a deck is better understood as a small domain-specific language (DSL): a simulator-specific program whose tags name GEOS components and whose attributes set their parameters. For example, a physics-module tag selects which governing equations are advanced, a material-model tag defines the constitutive law, and region or output tags connect those choices to parts of the mesh and requested results. Writing a deck is therefore a translation from natural-language intent into a structured program over the simulator's executable interface. This is difficult for four reasons: (i) the vocabulary is large and contains near-duplicate component names, such as \texttt{DruckerPrager}, \texttt{ExtendedDruckerPrager}, \texttt{DruckerPragerHardening}, and \texttt{ViscoDruckerPrager}; (ii) documentation and examples may reflect older interface versions; (iii) constraints span sections, so names introduced in one place must be reused consistently elsewhere; and (iv) task briefs often under-specify choices that experts fill in with domain-conventional defaults. The first difficulty is exact interface-token selection, while the others are program-level consistency constraints. A concrete annotated deck is in App.~\ref{app:geos-example}.


%% ============================================================
\section{Method}
\label{sec:method}

\begin{figure}[t]
\centering
% \includegraphics[width=\textwidth]{assets/siga_fig2.png}
\includegraphics[width=\textwidth]{assets/geos_fig2_v0.png}
\caption{\textbf{Execution trace of the SIGA agent loop.} The agent synthesizes natural language instructions, a system primer, and repository data to iteratively author an executable configuration deck. It cycles through gathering context via tool queries, taking action to draft the XML, and verifying results with a lint-checker until the deck is structurally valid and ready for simulation.}
\label{fig:2}
\end{figure}

% \subsection{Design-space components}
% \label{subsec:overview}

% Rather than re-implement yet another ReAct-style agent loop, we build on the substantial engineering already present in a frontier coding agent. Claude Code provides file access, shell, editing tools, retry logic, and an iterative coding workflow; we treat it as a base harness and use only its own extension mechanisms to inject domain grounding. We see this adaptation-over-reconstruction stance as a strength rather than a concession: it is a more repeatable recipe for new domains and target software, since each new simulator only requires writing a thin adapter rather than maintaining a full agent codebase; and the longer-term goal of fully automatic harness construction and optimization (\S\ref{subsec:self-evolved}) is more tractable when the optimization surface is a small adapter than a sprawling agent codebase that exposes many opportunities for over-specification. We introduce four binary components below, each motivated by a known agent failure mode rather than by a feature we wanted to ship. Implementation details are in App.~\ref{app:impl}; an execution trace is in Fig.~\ref{fig:2}.

% \textbf{R: retrieval over simulator artifacts.} The agent does not reliably know GEOS's vocabulary: class names, schema constraints, and example combinations of physics modules sit outside model pretraining. The bare harness can already explore the same GEOS documentation and example tree via its built-in filesystem tools (find/grep/Read); R provides an additional avenue, semantic search via embeddings, against the same knowledge base. R is a Model Context Protocol (MCP) server exposing three search tools: \texttt{search\_navigator} (GEOS Sphinx documentation pages), \texttt{search\_schema} (entries from the GEOS XSD schema), and \texttt{search\_technical} (example XML files and technical snippets). Each tool queries a separate ChromaDB collection built using OpenAI \texttt{text-embedding-3-small} embeddings. The failure mode addressed is unknown-vocabulary substitution.

% \textbf{S: stop-hook schema verification.} Self-refinement, in which the model is prompted to reflect on and improve its own output, is a well-known recipe for raising reliability, typically implemented as a control loop applied to the workflow's final answer. In the adapt-an-existing-harness setting, this control naturally takes the form of a termination hook. Even when the agent has authored a deck it believes complete, it sometimes terminates with an empty, unparseable, or schema-invalid file: a silent-incompleteness failure that downstream evaluation cannot rescue. S is a termination hook that intercepts every attempted termination, scans \texttt{/workspace/inputs/} for parseable XML, runs \texttt{xmllint --schema} against the canonical \texttt{.xsd}, and either allows termination or returns structured repair feedback. A per-task retry counter bounds re-prompts (default 2). Unlike agent-callable tools, this component is mandatory whenever enabled: every termination passes through it.

% \textbf{X: agent-callable XML validator.} While end-of-turn self-refinement is known to work, it confines validation to the workflow's endpoint. A natural extension is to give the agent the ability to actively check its own work mid-trajectory, at its discretion, rather than only hearing feedback at the end. X is the MCP tool \texttt{mcp\_\_xmllint\_\_validate\_geos\_xml}, exposing the same schema check as an optional tool the agent can choose to invoke. In enabled cells, the agent uses it about three times per task on average. The failure mode addressed is in-loop schema-violation drift.

% \textbf{M: procedural-memory cheatsheet.} A persistent memory scratchpad is a sample-efficient way to retain lessons and crucial domain knowledge from prior experience, sparing the agent from re-deriving the same vocabulary on every trajectory. M is a 775-token reference appended to the system prompt via \texttt{--append-system-prompt}: a compact vocabulary dump (physics-module families, XML element names, constitutive-model names, typical attributes, and short exemplars), distilled offline from 18 training trajectories using \texttt{google/gemini-3-flash-preview}, a different model from the inference model to avoid self-distillation. It is not an explicit policy. The failure mode addressed is recurring-vocabulary lookup.

% This four-component design space (R, S, X, M) is the object of study. We refer to any wrapper instantiating a subset of these components as a \textbf{Simulator-Interface Grounding Adapter (SIGA)}, with the corresponding factor letters denoting which components are enabled.
%%%%%%%%%%%%%%4.1

\subsection{Design-space components}
\label{subsec:overview}

We formulate simulator setup as an \emph{interface-grounding} problem for an off-the-shelf coding agent. 
The base harness already provides the generic agent loop: it can inspect files, edit candidate configurations, run shell commands, observe errors, and revise its outputs. 
What it lacks is not another ReAct-style control loop, but the executable contract of the target simulator: the simulator-specific vocabulary, structural constraints, validation rules, and conditions under which a generated workspace is complete enough to return. 
SIGA therefore takes an adaptation-over-reconstruction approach. 
We keep the coding harness fixed and inject a thin grounding layer through the harness's own extension mechanisms.

This design choice is intentional. 
A thin adapter is easier to port to new scientific software because only the simulator-facing contract must be rewritten, while the generic coding workflow remains unchanged. 
It also creates a smaller optimization surface for self-evolution (\S\ref{subsec:self-evolved}): prior trajectories can refine the adapter's primer, memory, and validation interfaces without modifying the entire agent loop. 
We instantiate this grounding layer as four binary components, each corresponding to a distinct failure point in simulator setup: vocabulary access, termination-time completeness, in-trajectory validation, and recurring procedural knowledge. 
Implementation details are in App.~\ref{app:impl}; an execution trace is shown in Fig.~\ref{fig:2}.

\paragraph{R: retrieval for simulator vocabulary grounding.}
The base agent may not reliably recover GEOS-specific tokens, such as solver names, material models, schema attributes, and example-compatible block combinations. 
Although the harness can search the repository with standard filesystem tools, keyword search is brittle when the agent does not already know the right simulator terms. 
R provides semantic access to the same simulator artifacts through a Model Context Protocol (MCP) server with three search tools: \texttt{search\_navigator} over GEOS documentation pages, \texttt{search\_schema} over XSD schema entries, and \texttt{search\_technical} over example XML files and technical snippets. 
This component targets unknown-vocabulary substitution.

\paragraph{S: stop-hook verification for termination validity.}
A frequent failure in simulator setup is silent incompleteness: the agent stops after producing a plausible-looking workspace, but the returned deck is empty, unparseable, schema-invalid, or missing required structure. 
S addresses this failure at the point where it matters most: termination. 
Whenever the agent attempts to finish, the stop hook scans \texttt{/workspace/inputs/} for parseable XML, runs \texttt{xmllint --schema} against the canonical GEOS \texttt{.xsd}, and either allows termination or returns structured repair feedback. 
A per-task retry counter bounds the number of re-prompts. 
Unlike an ordinary tool, S is mandatory whenever enabled: every attempted termination must pass through the verifier. 
This component targets invalid or incomplete final workspaces.

\paragraph{X: agent-callable validation for in-trajectory repair.}
Termination-time verification catches errors only at the end of the workflow. 
X instead exposes the same schema check as an optional MCP tool, \texttt{mcp\_\_xmllint\_\_validate\_geos\_xml}, that the agent may call while drafting the deck. 
This separates forced validation at termination from discretionary validation during construction: the agent can inspect an intermediate XML file, receive schema feedback, and revise before attempting to finish. 
In enabled cells, the agent invokes this validator about three times per task on average. 
This component targets schema-violation drift during the editing trajectory.

\paragraph{M: always-on procedural memory.}
Some simulator knowledge is repeatedly useful across tasks and should not have to be rediscovered from scratch on every trajectory. 
M provides this knowledge as an always-visible procedural-memory cheatsheet appended to the system prompt via \texttt{--append-system-prompt}. 
The cheatsheet is a compact 775-token reference containing GEOS physics-module families, XML element names, constitutive-model names, typical attributes, and short exemplars, distilled offline from 18 training trajectories using a model different from the inference model. 
It is not an explicit policy; it is a lightweight memory surface that keeps high-frequency simulator conventions in context. 
This component targets recurring-vocabulary lookup and repeated domain-interface mistakes.

Together, these four components define the SIGA design space. 
We refer to any wrapper instantiating a subset of $\{\mathrm{R,S,X,M}\}$ as a \textbf{Simulator-Interface Grounding Adapter (SIGA)}, with the corresponding factor letters denoting which grounding interventions are enabled.

%%%%%%%%%%%%%%%%%%%%%%%%%%

% \subsection{Self-evolved adapters}
% \label{subsec:self-evolved}

% The hand-designed components above were chosen from prior literature. A natural follow-up is whether an offline search over adapter contents, scored against held-out training-set TreeSim, can match or beat the best hand-designed cell. Treating the agent's own scaffolding as a learnable object (variously framed as meta-harness design, harness-as-code, agentic harness engineering, or skill optimization) is a rapidly emerging research direction, but most existing work targets core AI benchmarks (coding, terminal navigation, math). An open question is how this paradigm fares on a task driven by domain knowledge rather than by general programming competence. We provide an initial study of this direction in the SIGA setting. Our self-evolved variant is fundamentally similar to existing self-improvement and meta-harness methods in that the agent reflects on its own trajectories and writes its own plugin (the adapter); the hand-designed components above define the constrained design space within which this self-evolution operates. We evaluate two self-evolved variants:

% \textbf{SE: self-evolved combined adapter.} An offline pipeline iteratively rewrites the adapter contents, including a richer \texttt{PRIMER.md} and agent-authored skills, scoring revisions against held-out training-set TreeSim. At evaluation time we disable skill invocation with \texttt{--disallowedTools Skill} for tool-shape parity with the factorial cells.

% \textbf{SE-prose: prose-only self-evolved variant.} This variant takes only the rewritten v3 primer and cheatsheet from SE and inserts them into an otherwise standard S+X+M cell. It tests whether the SE gains come from improved system prompt and memory rather than from the full self-evolved adapter package. If SE-prose recovers most of the SE gain over S+X+M, the discovery process is operating on the prose channels.

%%%%%%%%%%%%%%%4.2

\subsection{Self-evolved adapter contents}
\label{subsec:self-evolved}

The four components above define where simulator-interface grounding can enter the coding harness, but their textual contents remain manually specified. 
We therefore test a modest form of adapter self-evolution: using prior trajectories to revise the lightweight grounding layer, then selecting revisions by validation-set TreeSim. 
This procedure does not modify the underlying coding agent, introduce a new planning policy, or search over the full agent loop. 
Instead, it optimizes the simulator-facing adapter contents, including the system primer, procedural memory, and auxiliary guidance, while keeping the base harness fixed.

This setting lets us ask whether the adapter can improve from its own execution history in a domain-knowledge-heavy task. 
Prior work on self-improving agents and harness optimization often treats the agent scaffold itself as a learnable object, but most demonstrations focus on general coding, terminal, or reasoning benchmarks. 
Here, the optimization target is narrower: the executable contract between a general coding agent and a scientific simulator. 
We evaluate two self-evolved variants to separate gains from improved prose grounding versus the broader self-evolved adapter package.

\paragraph{SE: self-evolved combined adapter.}
SE is produced by an offline pipeline that iteratively rewrites the adapter contents, including a richer \texttt{PRIMER.md} and agent-authored auxiliary skills, and scores candidate revisions by held-out training-set TreeSim. 
At evaluation time, we disable skill invocation with \texttt{--disallowedTools Skill} so that SE has the same tool shape as the factorial cells. 
Thus, any gain in the main comparison reflects the evolved adapter contents rather than an additional class of runtime tools.

\paragraph{SE-prose: prose-only self-evolved variant.}
SE-prose isolates the textual part of the self-evolved adapter. 
It takes only the rewritten primer and cheatsheet from SE and inserts them into an otherwise standard S+X+M cell. 
Comparing SE-prose with S+X+M tests whether self-evolution improves the always-visible grounding instructions and procedural memory. 
Comparing SE-prose with SE tests whether any remaining gain comes from the broader self-evolved adapter package beyond prose and memory.

%%%%%%%%%%%%%%%%%

%% ============================================================
\section{Experiments}
\label{sec:eval}

% \subsection{Research questions}
% \label{subsec:rqs}

% Our experiments are organised around four research questions that the results in \S\ref{sec:results} answer in turn.

% \textbf{RQ1 (mechanism).} Which grounding components carry the empirical weight, and on which axis (reliability vs.\ absolute quality) do they help? Answered in \S\ref{subsec:quality} and \S\ref{subsec:bottleneck-results}.

% \textbf{RQ2 (transfer).} Does the same grounding recipe transfer beyond GEOS to a different simulator family with a different interface? Answered in \S\ref{subsec:openfoam-transfer}.

% \textbf{RQ3 (autonomy).} When the brief specifies less, how does the agent behave, and does it consult a human supervisor to resolve the ambiguity? Answered in \S\ref{subsec:autonomy}.

% \textbf{RQ4 (human anchor).} How does the SIGA-equipped agent compare with domain experts on a representative task? Answered in \S\ref{subsec:human-baseline}.

%%%%%%%%%%%%% 5.1
\subsection{Experimental Design}

Our experiments are designed to separate three questions that are often conflated in simulator-agent evaluation: 
whether the adapter improves final deck quality, whether it improves reliability by preventing catastrophic invalid 
outputs, and which simulator-interface failures remain unsolved. We first run a controlled factorial ablation on GEOS 
to estimate the effect of each grounding component. We then analyze failure categories to identify the residual 
bottlenecks. Finally, we calibrate the task against human domain users, test behavior under underspecified briefs, 
and run a small OpenFOAM transfer study to probe whether the dominant mechanism is GEOS-specific.

The results answer five research questions:
\textbf{RQ1: Reliability and quality.} Does SIGA improve TreeSim, and is the gain driven by average quality or by fewer catastrophic failures?
\textbf{RQ2: Failure mechanisms.} Which block-level or attribute-level errors are fixed by the adapter?
\textbf{RQ3: Human anchor.} How does SIGA compare with domain experts on a representative GEOS deck-authoring task?
\textbf{RQ4: Autonomy.} When task briefs are underspecified, does the agent consult a human supervisor or substitute other available information sources?
\textbf{RQ5: Transfer.} Does the same grounding mechanism transfer to a different simulator interface?

%%%%%%%%%%%%%%%

%%%%%%%%%%%%%%%%% 5.2

% \subsection{Controlled ablation over grounding components}

% To attribute gains to adapter mechanisms rather than to a single monolithic wrapper, we evaluate SIGA as a controlled 
% design space over four binary factors: retrieval (R), stop-hook validation (S), agent-callable validation (X), and 
% procedural memory (M). A full $2^4$ factorial would require 16 cells. Because our goal is to estimate main effects 
% while controlling compute, we use a Resolution-IV $2^{4-1}$ fractional factorial design. This preserves unconfounded 
% main-effect estimates while requiring only eight factorial cells. We additionally evaluate S+X+M, SE-prose, and SE 
% because they correspond to the predicted best non-retrieval adapter, the prose-only self-evolved adapter, and the full 
% self-evolved adapter.

%%%%%%%%%%%%%%%%%

\subsection{Ablation design and bottleneck analysis}
\label{subsec:cells}

\textbf{Cells.} To attribute the contribution of each grounding component, we run a detailed factorial ablation across the four binary factors $\{\mathrm{R, S, X, M}\}$. A standard one-factor-at-a-time ablation (stacking each component on, or stripping it from the full stack) is the cheapest design, but it cannot disambiguate main effects from two-factor interactions: if R and S help only when combined, neither single-factor ablation reveals this. The full $2^4 = 16$-cell factorial does, at the cost of $4{\times}$ the runs. We instead use a \emph{Resolution-IV $2^{4-1}$ fraction} with generator $\mathrm{D}{=}\mathrm{ABC}$, which gives us eight cells whose main effects are not confounded with two-factor interactions, recovering most of the information of the full factorial at half the compute. Cells are named by their active factors, e.g., $\mathrm{X+M}$; \emph{Vanilla} is all-off. We add three cells: $\mathrm{S+X+M}$, the predicted main-effects-best corner missing from the fraction; \emph{SE-prose}, which inserts the self-evolved v3 primer and cheatsheet into an otherwise standard $\mathrm{S+X+M}$ cell; and \emph{SE}, the full self-evolved combined adapter. Full settings are in App.~\ref{app:cells}. Because S and X both use \texttt{xmllint}, the X main effect partly conflates agent-callable validation with hook-time schema validation when S is also enabled (\S\ref{sec:discussion}).

\textbf{Bottleneck analysis.} We attribute score differences to structural failure modes in three steps (schema and prompts in App.~\ref{app:bottleneck}). First, an LLM-free extractor identifies top-$K$ failing subtrees from \texttt{treesim\_detail} and adds trajectory features from \texttt{events.jsonl}. Second, \texttt{deepseek-v4-flash} assigns each (cell, seed, task) to one category: \texttt{missing\_block}, \texttt{extra\_block}, \texttt{hallucinated\_extras}, \texttt{structural\_mismatch}, \texttt{bad\_attribute\_value}, \texttt{partial\_implementation}, \texttt{wrong\_constitutive}, or \texttt{no\_failure}. Third, \texttt{deepseek-v4-pro} summarizes per-cell distributions and baseline-vs-best deltas. Total cost is ${\sim}650$ flash calls plus 4 pro calls ($\$5\text{--}7$).

\textbf{Splits and leakage control.} The 46 tasks are split into 10 held-out-eval tasks, 18 distillation tasks, and 17 validation-selection tasks (App.~\ref{app:bench}). No validation or held-out-eval task is used for distillation. A hygiene gate regex-scans distilled artifacts for ground-truth test filenames and rejects leaks. (It was added after an earlier cheatsheet leaked 13/17 validation-task basenames.)

% \subsection{Benchmark}
% \label{subsec:benchmark}
% We evaluate on two task sets (App.~\ref{app:bench}). \textbf{val} contains 17 in-distribution GEOS advanced-example and tutorial tasks spanning poromechanics, hydraulic fracture, thermal coupling, and wellbore modeling (e.g., \texttt{ExampleDPWellbore}, \texttt{ExampleMandel}, \texttt{buckleyLeverettProblem}). These form the held-out half of a 36-task index-parity split; cells were selected against this set, the cheatsheet was not. \textbf{held-out-eval} contains 10 harder tasks reserved for final evaluation, never used for distillation, self-evolution, or cell tuning. Inputs are v2 natural-language specifications; ground-truth decks are hand-validated.

% \subsection{Metric: TreeSim}
% \label{subsec:metric}
% We score generated decks with \textbf{TreeSim}, a tree-edit similarity in $[0,1]$ decomposed by the ten canonical GEOS sections. Headline numbers average TreeSim under \emph{failures-as-zero}: parse errors, timeouts, \texttt{failed\_no\_outputs}, and missing XML outputs all score 0, so systems are not rewarded for unscorable files.

% \subsection{Baselines, models, and runs}
% \label{subsec:baselines}
% The main baseline is Vanilla Claude Code with all SIGA components disabled. \emph{SE} and \emph{SE-prose} use self-evolved adapter content (\S\ref{subsec:self-evolved}). A one-shot LLM baseline gives a lower-bound reference (App.~\ref{app:cross-model-detail}; $n=1$, minimax-m2.7, val TreeSim $=0.333$). The OpenFOAM transfer (\S\ref{subsec:openfoam-transfer}) compares against Foam-Agent 2.0~\citep{foamagent2025} restricted to its linting tool. All headline GEOS runs use \texttt{deepseek-v4-flash}, $n=3$ seeds per cell, default Claude Code sampling, 1500~s per-task timeout. Bottleneck classification uses \texttt{deepseek-v4-flash}; cross-cell synthesis uses \texttt{deepseek-v4-pro}. SE cells run with \texttt{--disallowedTools Skill} for parity.


%%%%%%%%%%%%5.3 - 5.5
\subsection{GEOS benchmark and evaluation protocol}

We evaluate on two GEOS task splits. The validation split contains 17 in-distribution tasks from advanced examples 
and tutorials, covering poromechanics, hydraulic fracture, thermal coupling, and wellbore modeling. This split is used 
for component selection but not for memory distillation. The held-out-eval split contains 10 harder tasks reserved 
for final evaluation and is never used for distillation, self-evolution, or cell tuning.

Generated decks are scored with TreeSim, a tree-edit similarity metric in $[0,1]$ computed over the canonical GEOS 
deck structure. We report failures-as-zero: parse errors, timeouts, missing XML outputs, and empty outputs all receive 
score 0. This convention is important because simulator setup is only useful when the returned workspace is executable 
or at least structurally inspectable.

The main baseline is vanilla Claude Code with all SIGA components disabled. All headline GEOS experiments use 
\texttt{deepseek-v4-flash} with three seeds per cell and a 1500-second timeout. SE and SE-prose use the same evaluation budget 
and disable skill invocation for parity.

%%%%%%%%%%%%


%% ============================================================
\section{Results}
\label{sec:results}

The main empirical finding is that simulator-interface grounding improves reliability more than it improves the 
best-case quality ceiling. On in-distribution GEOS tasks, vanilla Claude Code already achieves high TreeSim, leaving 
little room for uniform average-quality gains. On harder held-out tasks, however, SIGA reduces catastrophic invalid or 
incomplete outputs, which raises mean performance and sharply reduces across-seed variance. The remaining errors are 
not mostly structural omissions but attribute-level semantic mistakes, suggesting that current adapters raise the floor 
rather than solve simulator reasoning.

% \subsection{Reliability, not absolute quality, is where adapters help (RQ1)}
\subsection{SIGA raises the reliability floor rather than the quality ceiling (RQ1)}
\label{subsec:quality}

Table~\ref{tab:main-results} reports headline numbers ($n=3$). Cells cluster narrowly on val (0.857 to 0.921; 0.910 to 0.921 for non-RAG) but spread more on held-out-eval (0.720 to 0.789, $\Delta = +0.069$ Vanilla to SE). The gap between these two columns, more than the absolute level of either, is the headline finding: the adapters' contribution is largely reliability rather than uniform quality lift.


\begin{table}[t]
  \caption{Cell-level TreeSim (failures-as-zero) on \texttt{deepseek-v4-flash}, $n=3$ seeds (mean $\pm$ sample std). The left block indicates which grounding components are present in each cell (\textbf{$\bullet$} on, blank off): \emph{Retrieval} (semantic search over GEOS artifacts), \emph{Refine loop} (schema-validating termination hook), \emph{Validator} (the same schema check exposed as an agent-callable tool), and \emph{Memory} (procedural-memory cheatsheet appended to the system prompt). The first eight rows are a Resolution-IV $2^{4-1}$ fractional factorial; \emph{S+X+M} fills the main-effects-best corner missing from the fraction; \emph{SE-prose} substitutes the self-evolved primer and cheatsheet into an otherwise standard \emph{S+X+M} cell; \emph{SE} is the full self-evolved adapter. $\Delta$ = cell mean minus \emph{Vanilla} mean on the same split.}
  \label{tab:main-results}
  \centering
  \small
  \resizebox{0.98\textwidth}{!}{%
  \begin{tabular}{l|cccc|cccc}
    \toprule
    & \multicolumn{4}{c|}{\textbf{Components}} & \multicolumn{4}{c}{\textbf{TreeSim}} \\
    \textbf{Cell} & \textbf{Retrieval} & \textbf{Refine loop} & \textbf{Validator} & \textbf{Memory} & \textbf{val} & \textbf{held-out-eval} & \textbf{$\Delta$ val} & \textbf{$\Delta$ held-out} \\
    \midrule
    Vanilla        &         &         &         &         & $0.910 \pm 0.024$           & $0.720 \pm 0.081$            & --                & --                \\
    R$+$M          & $\bullet$ &         &         & $\bullet$ & $0.885 \pm 0.014$           & --                           & $-0.025$          & --                \\
    S$+$M          &         & $\bullet$ &         & $\bullet$ & $0.919 \pm 0.004$           & --                           & $+0.009$          & --                \\
    R$+$S          & $\bullet$ & $\bullet$ &         &         & $0.857 \pm 0.045$           & --                           & $-0.053$          & --                \\
    X$+$M          &         &         & $\bullet$ & $\bullet$ & $\mathbf{0.921 \pm 0.007}$  & $0.768 \pm 0.005$            & $\mathbf{+0.011}$ & $+0.048$          \\
    R$+$X          & $\bullet$ &         & $\bullet$ &         & $0.893 \pm 0.033$           & --                           & $-0.017$          & --                \\
    S$+$X          &         & $\bullet$ & $\bullet$ &         & $0.917 \pm 0.004$           & $0.781 \pm \mathbf{0.002}$   & $+0.007$          & $+0.061$          \\
    R$+$S$+$X$+$M  & $\bullet$ & $\bullet$ & $\bullet$ & $\bullet$ & $0.885 \pm 0.008$           & --                           & $-0.025$          & --                \\
    \midrule
    S$+$X$+$M      &         & $\bullet$ & $\bullet$ & $\bullet$ & $0.911 \pm 0.018$           & $0.783 \pm 0.022$            & $+0.001$          & $+0.063$          \\
    SE-prose       &         & $\bullet$ & $\bullet$ & $\bullet$ & $0.897 \pm 0.032$           & $0.775 \pm 0.024$            & $-0.013$          & $+0.055$          \\
    SE             &         & $\bullet$ & $\bullet$ & $\bullet$ & $0.919 \pm 0.020$           & $\mathbf{0.789 \pm 0.012}$   & $+0.009$          & $\mathbf{+0.069}$ \\
    \bottomrule
  \end{tabular}}
\end{table}

\textbf{Hard-tail rescue.} Concretely, on held-out-eval Vanilla has $\sigma = 0.081$ because a single seed produced unparseable XML on \texttt{ExampleProppantTest} and scored 0; under X+M this drops to $\sigma = 0.005$ and under S+X to $\sigma = 0.002$, roughly an order-of-magnitude reduction in across-seed standard deviation. The underlying mechanism is the prevention of a small number of zero-score outputs on the harder task set: the adapter's first-order job is to keep the harness from terminating with empty or unparseable decks, and mean lift on the harder split follows from rescuing those would-be zero-score runs. Per-task inspection (App.~\ref{app:per-task}, \S\ref{subsec:per-task}) shows the $+0.069$ Vanilla-to-SE aggregate gain on held-out-eval is driven primarily by two catastrophic-failure rescues (\texttt{AdvancedExampleThermoPoroElasticWellbore} from $0.355$ to $0.761$; \texttt{ExampleProppantTest} from $0.541$ to $0.825$), alongside one universal failure (\texttt{TutorialHydraulicFractureWithAdvancedXML} at $0.013$ across every cell) and within-noise differences on the remaining tasks. The adapter changes rescue the bare model from these catastrophic failures rather than uniformly raising scores: on tasks where the bare model already has a usable template, the adapter is operating in seed noise.

\label{subsec:per-task}\textbf{Val ceiling.} On val, where the bare harness already operates near a quality ceiling, the absolute-score improvements from any single component are subtle. The best val cell (X+M) beats Vanilla by $+0.011$, within seed std, and we are honest that this is a modest absolute lift; the headline contribution on val is again on the reliability axis (variance across seeds) and on efficiency (\S\ref{subsec:efficiency}) rather than absolute quality. The only Resolution-IV main effect on val that clears the noise floor is R, with a small negative sign ($\Delta = -0.032$); X, M, and S all fall within $\pm 0.007$ (App.~\ref{app:main-effects}). 
% We flag that the R main-effect estimate is partly contaminated by the native-plugin-prefix bug described in App.~\ref{app:cross-model-detail}: cells with $\mathrm{R^-}$ that did not opt out of the prefix were instructed to call retrieval tools that were not registered, so the $\mathrm{R^-}$ side of the contrast underweighted that subset. The X+M-vs-Vanilla contrast is unaffected (both are $\mathrm{R^-}$ in the same way), but the reported $-0.032$ R effect should be treated as an upper bound on R's negative magnitude until a clean re-run completes (App.~\ref{app:future-work}).

\textbf{Takeaway (RQ1).} Adapters move the reliability axis, not the absolute-quality ceiling: they shrink across-seed variance by roughly an order of magnitude on held-out-eval and rescue two compound multi-physics tasks from catastrophic failure, but they do not raise the count of strictly-correct decks on the in-distribution split.

\subsection{Schema-aware adapters fix block-level omissions; attribute errors persist (RQ2)}
\label{subsec:bottleneck-results}

The per-task classifier (App.~\ref{app:bottleneck-counts}) reveals four patterns. \textbf{(1) Adapters fix \texttt{missing\_block}:} counts drop 6 to 3 on val between Vanilla and X+M (whole \texttt{Solvers}, \texttt{Events}, or \texttt{Constitutive} omissions), since schema validation requires these blocks and the cheatsheet enumerates them. \textbf{(2) Adapters do not fix \texttt{bad\_attribute\_value}} (12 / 11 / 15 across Vanilla / X+M / S+X): schema validation cannot prevent \texttt{TPFAstabilization} substituted for \texttt{contactStabilization} or a hallucinated \texttt{gravityVector}. \textbf{(3) Adapters trade catastrophic absence for content imprecision:} as \texttt{missing\_block} drops, \texttt{extra\_block} rises 9 to 11 and \texttt{hallucinated\_extras} rises 4 to 7, yielding small mean lift but large reliability lift. \textbf{(4) Strictly perfect tasks (TreeSim $\geq 0.999$) do not increase} under any adapter (Vanilla 7/51, X+M 6/51, SE 6/51). The harness operates in a harm-reduction regime, not a correctness regime.

\paragraph{Efficiency.}\label{subsec:efficiency} Adapter cells are no more expensive than Vanilla in wall-clock (App.~\ref{app:results}). This matched-efficiency outcome is itself a feature rather than a null finding. A common failure mode in LLM-reflection-driven harness optimization is over-specification, where iteratively tacked-on small features end up derailing or slowing down workflows, so a non-trivial design constraint is that the adapter not impose runtime overhead beyond the bare harness. Our adapters meet this constraint: cheatsheet cells make about half as many free-form \texttt{Read} calls because the cheatsheet short-circuits exploratory file access; xmllint cells make roughly 3 schema-validation calls per task; RAG cells make 12 to 13 retrieval calls and run faster but score lower. SE makes fewer tool calls than Vanilla on val (68.9 vs 81.5, ${-}15.5\%$) but more on held-out-eval (97.4 vs 90.5, $+7.6\%$); the val efficiency does not transfer to the harder task split, consistent with the broader pattern that adapter wins are tail-localised.

\textbf{Takeaway (RQ2).} Schema-aware adapters reliably fix block-level omissions but leave attribute-level semantic errors untouched, exposing the residual frontier; cells also trade catastrophic absence for content imprecision, which yields the small mean lift and large reliability lift we observe. Adapters add no wall-clock overhead.

\subsection{Human baseline: SIGA reaches expert deck quality on a representative task (RQ3)}
\label{subsec:human-baseline}

To anchor agent TreeSim in human terms, two geoscience-domain-expert volunteers (Expert 1, Expert 2; graduate level geoscientists, new to GEOS) attempted \texttt{buckleyLeverettProblem} (1D immiscible CO$_2$/brine displacement, the easy end of our bench) within a one-hour budget, using the same primer, filtered source tree, and contamination block as the agent, and no LLM chatbots or web search. Expert 1 returned with no time cap as an extended-budget check, and we collected a written estimate from a GEOS expert and developer to anchor experienced users (App.~\ref{app:human-browser}).

Both 1h participants ran out of time at 47 to 48 min and submitted only \texttt{base.xml}, with the second half spent on \texttt{Outputs} and \texttt{Events} blocks. File-level TreeSim is $0.812$ (Expert 1) and $0.781$ (Expert 2); deck-level collapses to $0.540$ and $0.527$ since the second file is missing (Table~\ref{tab:human-baseline}). On the same task, vanilla Claude Code reaches file-level $0.889 \pm 0.023$ and deck-level $0.751 \pm 0.016$ in roughly 7 min; SIGA X+M reaches at least $0.90$ on both. The 1h shortfall reflects the budget, not capacity: Expert 1 with no cap produced both files in about 3 hours and reached deck-level $\mathbf{0.931}$, at parity with SIGA X+M but at roughly $36\times$ the wall-clock. The GEOS developer estimate brackets an experienced GEOS user at under 30 min for simple Buckley--Leverett and ``a couple of days'' for compound multi-physics decks, lining up with the hard-tail held-out-eval result.

\begin{table}[h]
  \caption{Human baseline on \texttt{buckleyLeverettProblem}. File-level TreeSim is GT \texttt{base.xml} vs the participant's \texttt{base.xml}; deck-level is the full GT directory vs the submitted directory. Agent rows from per-task \texttt{deepseek-v4-flash} numbers.}
  \label{tab:human-baseline}
  \centering
  \small
  \begin{tabular}{lrrr}
    \toprule
    \textbf{Author} & \textbf{File-level (base)} & \textbf{Deck-level} & \textbf{Wall (min)} \\
    \midrule
    Expert 1 (1\,h cutoff)                          & $0.812$            & $0.540$            & $48.2$ \\
    Expert 2 (1\,h cutoff)                          & $0.781$            & $0.527$            & $46.7$ \\
    Expert 1 (no time cap, both files)              & $0.689$            & $0.931$            & ${\sim}180$ \\
    Vanilla CC (\texttt{deepseek-v4-flash})   & $0.889 \pm 0.023$  & $0.751 \pm 0.016$  & ${\approx}\,7$ \\
    SIGA X+M (\texttt{deepseek-v4-flash})     & ${\geq}\,0.90$     & ${\geq}\,0.90$     & ${\approx}\,5$ \\
    \bottomrule
  \end{tabular}
\end{table}

Browser histories show humans browsing Sphinx prose to assemble decks from concept descriptions, while the agent analogises from prior decks via the source tree (App.~\ref{app:human-browser}). Part of the SIGA contribution is giving the agent a structured way to do what humans approximate via Sphinx browsing under time pressure. We treat this as an existence-of-effect calibration, not a head-to-head ranking.

\textbf{Takeaway (RQ3).} The SIGA-equipped agent reaches the deck-level quality of an extended-budget domain expert on this representative task at roughly $1/36$ the wall-clock, and operates inside the experienced-GEOS-user time envelope for simple Buckley--Leverett deck authoring.

\subsection{Exploring agent autonomy: the example library substitutes for human consultation (RQ4)}
\label{subsec:autonomy}

\textbf{Motivation.} Our headline benchmark hands the agent a fully detailed task brief in a single turn, by design: this standardised setup removes potential ambiguity and confounders, so any difference between adapter cells is attributable to the adapter rather than to varying levels of underspecification. In this regime the agent's responsibility is narrowed to translation: it is told exactly which simulation to run, and only needs to express it in GEOS's DSL. Eventually, however, we want to expand the scope of the agent's responsibility so that it can take on more of the workflow and be correspondingly more useful to a working scientist. Loosening the brief naturally introduces ambiguity that a human scientist collaborator would otherwise resolve. We explore this direction by producing alternative specifications that specify less and less of the task, requiring the agent to figure out more on its own, and observing when and how the agent chooses to consult a human expert user to resolve the ambiguities it encounters.

\textbf{Protocol.} We probe this with a single-seed companion study (8 val tasks, 2 configurations $\times$ 2 relaxation levels $\times$ 2 interaction modes, 64 runs; App.~\ref{app:autonomy}). Briefs are tier-rewritten by \texttt{deepseek-v4-pro}: \textbf{Medium} drops T1 (software defaults) and T2 (standard numerics), 89 values ($-21\%$ chars); \textbf{Hard} additionally drops T3 (domain-inferable physicals), 184 values ($-50\%$). Configurations are Vanilla and X+M; the \emph{interactive} mode exposes a \texttt{consult\_supervisor} MCP tool whose handler invokes a separate \texttt{deepseek-v4-flash} prompted with the full original brief.

\textbf{Results.} Across 64 interactive trials the agent invoked the exposed human-consultation tool only twice ($3.1\%$), a rate that is robust to prompt framing. The 31/32 silent runs were not silent for lack of questions: agents read 142 to 404 GEOS example XMLs per cell and copied values from analogous benchmarks. A neutral-framing rerun (no ``prefer to infer'' language) leaves the rate unchanged at 1/32, ruling out prompt-induced reluctance. A Mandel/Hard diagnostic shows 15 of 26 dropped values appear by literal-token \texttt{grep} in other readable GEOS examples, so the on-disk library is acting as a cheaper retrieval substitute for human consultation. TreeSim drops mildly with relaxation (X+M: $0.829$ Medium, $0.835$ Hard, vs $0.921$ at Easy), and interactive variants are within $\pm 1$pp of non-interactive: the agent reaches the same place whether or not the human channel exists.

\textbf{Takeaway (RQ4).} ``Decide when to consult human oversight'' is not a free-standing agent capability in this setting; it is conditional on whether the agent has access to an alternative oracle. The on-disk example library acts as a cheaper retrieval substitute for human consultation, so studies measuring consultation behaviour without controlling for executable-example-library access are measuring the substitute, not the consultation.
% We recommend physics-family-level contamination blocks (e.g., blocking all \texttt{*Mandel*} and \texttt{*Poroelastic*} input files for a Mandel-class task, not just the GT basenames) as a follow-up.

\subsection{Transfer pilot: forced end-of-turn verification also improves OpenFOAM coverage (RQ5)}
\label{subsec:openfoam-transfer}

% The 5-task OpenFOAM setting reported below is a pilot study, intended only to sanity-check that the SIGA methodology is not specific to GEOS rather than to serve as a second full benchmark. We ported $\{\mathrm{R,S,X,M}\}$ to OpenFOAM case authoring on a 5-task FoamGPT-derived subset (single seed, file-text-and-coverage metric, App.~\ref{app:openfoam-transfer}). The qualitative lesson transfers: the best cell (R+S) reaches mean $0.871$, versus $0.466$ for vanilla Claude Code and $0.569$ for Foam-Agent~\citep{foamagent2025} in lint-only mode (the stable mode in our environment). Every $\mathrm{S}$-enabled cell achieves full required-file coverage with no zero-score failures; Vanilla covers 3/5 and R+X covers 1/5. A factor-style readout assigns the largest effect to $\mathrm{S}$ ($+0.328$ mean), with $\mathrm{M}$ positive and $\mathrm{R}$, $\mathrm{X}$ slightly negative.

% \textbf{Takeaway (RQ2).} The dominant mechanism transfers: the same end-of-turn verification component (the refine loop) that carries the GEOS reliability lift also carries the OpenFOAM coverage lift, despite the different interface format and metric. We treat this as transfer evidence, not a second benchmark.


We next ask whether the dominant SIGA mechanism is specific to GEOS's XML/XSD interface or reflects a broader adapter principle: forcing the agent to satisfy simulator-specific completion checks before termination. 
To probe this question, we run a small OpenFOAM transfer pilot rather than a second full benchmark. 
We port $\{\mathrm{R,S,X,M}\}$ to OpenFOAM case authoring on a 5-task FoamGPT-derived subset, using a single seed and a file-text-and-coverage metric (App.~\ref{app:openfoam-transfer}). 
Because the task set is small and the metric differs from GEOS TreeSim, we use this study only as qualitative transfer evidence.

The qualitative pattern matches the GEOS reliability result. 
The best OpenFOAM cell, R+S, reaches a mean score of $0.871$, compared with $0.466$ for vanilla Claude Code and $0.569$ for Foam-Agent~\citep{foamagent2025} in lint-only mode, the stable mode available in our environment. 
More importantly, the gain is again driven by termination-time completeness enforcement: every $\mathrm{S}$-enabled cell produces all required files and avoids zero-score failures, whereas Vanilla covers only 3/5 tasks and R+X covers only 1/5. 
A factor-style readout assigns the largest positive effect to $\mathrm{S}$ ($+0.328$ mean), while the other factors are smaller and less stable in this pilot.

\textbf{Takeaway (RQ5).} 
The OpenFOAM pilot supports the same mechanism-level conclusion as GEOS: end-of-turn verification raises the reliability floor by preventing incomplete simulator workspaces from being returned. 
This is not evidence that SIGA is a universally best OpenFOAM agent, nor a full head-to-head benchmark against Foam-Agent's intended execution-coupled workflow. 
Rather, it suggests that forced termination checks are a transferable adapter principle across simulator interfaces: XSD validity in GEOS and required-file/dictionary coverage in OpenFOAM.


%% ============================================================
\section{Discussion}
\label{sec:discussion}
\label{sec:analysis}
\label{subsec:limitations}

\paragraph{What transfers across simulators.} Of the four components, only $\mathrm{S}$ shows a large positive effect on both GEOS reliability (roughly an order-of-magnitude $\sigma$ reduction, hard-tail rescue) and OpenFOAM coverage ($+0.328$ mean, full required-file coverage on every $\mathrm{S}$-enabled cell). $\mathrm{R}$ is negative on GEOS val and slightly negative on OpenFOAM; $\mathrm{X}$ is small-positive on GEOS and small-negative on OpenFOAM; $\mathrm{M}$ is small-positive on both. The transferable mechanism is forced end-of-turn verification, and a paper reading our results that wants to ship a single adapter component is best advised to ship $\mathrm{S}$.

\paragraph{New tools are not used just because they are exposed.} An earlier version of the adapter exposed cheatsheet content through an MCP \texttt{memory\_lookup} tool with embedding top-$k$ retrieval over the distilled corpus, motivated by the retrievable-procedural-memory framing common in self-evolving-agent literature. Across every test-set run in which the tool was available, the agent called it zero times (the tool was verified functional). Delivering the same content through always-on \texttt{--append-system-prompt} (the M factor) is what produced any lift. The general lesson is that providing the agent with a new tool does not guarantee the agent will choose to use it: a retrieval-based memory module depends on agent-initiated lookup, and should be benchmarked against an always-on alternative before being treated as effective.

\paragraph{Adapter-design recommendations.} Each is grounded in a failure category that survives the Vanilla-to-best-config transition. \textbf{(i)} Schema- or coverage-enforced block presence is the cheapest reliable adapter; the OpenFOAM transfer suggests this generalises to any end-of-turn completeness check, not just XSD validation. \textbf{(ii)} Procedural memory should be delivered as an always-on system-prompt augmentation rather than as a retrievable tool; an agent that has to decide to invoke memory may simply never do so. \textbf{(iii)} Memory cheatsheets enumerating ``for physics $X$, use solver $Y$'' should be paired with explicit negative constraints (``the GT contains exactly $k$ \texttt{Constitutive} children, no more''); without them they trade \texttt{missing\_block} for \texttt{extra\_block} and \texttt{hallucinated\_extras}. \textbf{(iv)} Attribute-level oracles (per-solver allowed-attribute tables) are the next frontier; \texttt{bad\_attribute\_value} is untouched by everything we tested. \textbf{(v)} Autonomy benchmarks should remove the easy on-disk oracles before measuring consultation behaviour. \textbf{(vi)} Closed-loop retries driven by validator output are needed to raise the ceiling, since static hooks only raise the floor.

% \paragraph{Limitations.} The $\mathrm{X}$ main effect conflates ``MCP validator alone'' with ``hook also runs \texttt{xmllint} when $\mathrm{S}$ is on.'' Headline numbers are $n=3$ on \texttt{deepseek-v4-flash}; cross-model results on \texttt{minimax-m2.7} and \texttt{gemini-3-flash-preview} are single-seed, and we did not run OpenCode. The OpenFOAM study is 5 tasks, $n=1$, with Foam-Agent in lint-only mode (execute-mode failed in our environment); a clean head-to-head in execute mode is the natural follow-up. TreeSim is structural, not physical: a $0.8$ deck is not guaranteed to run. The native-plugin-prefix bug (App.~\ref{app:cross-model-detail}) contaminated the $\mathrm{R}=-0.033$ estimate; the X+M-vs-Vanilla ranking is unaffected since both are $\mathrm{R^-}$. The held-out-eval lift is concentrated in two tasks; the broader claim ``adapters help on hard compound multi-physics tasks'' is supported by these two but a larger held-out set would tighten it. The human baseline is $n=2$ on a single task and the budget-matched comparison is bounded below by the agent: both humans timed out before producing both required files. We treat the wall-clock anchor as existence-of-effect, not a head-to-head ranking.

\paragraph{Broader impact.} Lowering the configuration barrier could accelerate legitimate research (CO$_2$ storage, geothermal, induced seismicity) but could also enable less careful studies. Human expert review remains important; we frame the system as an assistant.


%% ============================================================
\section{Conclusion}
\label{sec:conclusion}

We asked how far an off-the-shelf coding agent can be pushed for scientific simulator setup by wrapper-level grounding alone, taking GEOS deck authoring as our case study. On in-distribution tasks the bare harness already operates near a quality ceiling; on harder held-out compound multi-physics tasks our adapters lift mean TreeSim by roughly seven percentage points and reduce across-seed variance by roughly an order of magnitude, largely by preventing the harness from occasionally returning unparseable or empty decks. Schema-driven adapters eliminate whole-block omissions but leave attribute-level errors untouched, and an OpenFOAM pilot reproduces the same dominant mechanism. As an application, SIGA brings deck authoring inside the time envelope of an experienced human user; as an initial AI-for-science contribution, it argues that the advanced tool-operating bottleneck deserves deliberate attention before agents are entrusted with harder scientific reasoning, and that adapting an existing engineered harness is a viable alternative to rebuilding agents from scratch for each new scientific target.
% Acknowledgments hidden for anonymous submission; restore for camera-ready.

\medskip

{
\small
}
\bibliographystyle{abbrvnat}
\bibliography{references}


%%%%%%%%%%%%%%%%%%%%%%%%%%%%%%%%%%%%%%%%%%%%%%%%%%%%%%%%%%%%

\appendix

\section{Benchmark details}
\label{app:bench}

\subsection{Task pool, splits, and hygiene}
The 46-task pool is mined from GEOS advanced examples and tutorial decks. Ten tasks are held out as an ICL pool: \texttt{AdvancedExampleCasedThermoElasticWellbore}, \texttt{AdvancedExamplePureThermalDiffusionWellbore}, \texttt{AdvancedExampleThermoPoroElasticWellbore}, \texttt{AdvancedExampleViscoExtendedDruckerPrager}, \texttt{ExampleIsothermalHystInjection}, \texttt{ExampleMCCWellbore}, \texttt{ExampleProppantTest}, \texttt{ExamplesingleFracCompression}, \texttt{ExampleVerticalPoroElastoPlasticWellbore}, \texttt{TutorialHydraulicFractureWithAdvancedXML}. The remaining 36 are split by odd/even index (sorted by training-run TreeSim, failures-as-zero) into an 18-task distillation corpus and a 17-task test set (one task dropped). Ground-truth decks live at \texttt{/data/shared/geophysics\_agent\_data/data/eval/experiments\_gt/}.

\subsection{Spec versions (v1 vs v2)}
Two instruction sets exist: v1 (original mined specs, used in pre-2026-04-21 deepseek runs) and v2 (cleaner, more self-contained rewrites, used in all post-D-004 runs). Mid-project comparisons that crossed the versions were contaminated; re-runs on v2 re-anchored the canonical baselines. All numbers in this paper use v2.

\section{Cell definitions}
\label{app:cells}

\begin{table}[h]
  \caption{Cell definitions. \emph{Retrieval}=semantic search over GEOS artifacts (RAG plugin); \emph{Refine loop}=schema-validating termination hook (parse-check; uses xmllint when \emph{Validator} is also on); \emph{Validator}=schema check exposed as an agent-callable tool (xmllint MCP); \emph{Memory}=procedural-memory cheatsheet. The first eight rows are the Resolution-IV $2^{4-1}$ factorial; \emph{S+X+M} fills in the missing 16th cell at the predicted main-effects-best corner; \emph{SE-prose} substitutes the v3 prose primer and cheatsheet for m1u while keeping \emph{S+X+M} factor settings; \emph{SE} is the self-evolved monolithic plugin (skills disabled at evaluation time).}
  \label{tab:cells}
  \centering
  \small
  \begin{tabular}{lccccl}
    \toprule
    \textbf{Cell} & \textbf{Retrieval} & \textbf{Refine loop} & \textbf{Validator} & \textbf{Memory} & \textbf{Notes} \\
    \midrule
    Vanilla &         &         &         &         & vanilla Claude Code (factorial baseline)               \\
    R+M & $\bullet$ &         &         & $\bullet$ & retrieval + memory                                     \\
    S+M &         & $\bullet$ &         & $\bullet$ & refine loop (parse-check) + memory                     \\
    R+S & $\bullet$ & $\bullet$ &         &         & retrieval + refine loop (parse-check)                  \\
    X+M &         &         & $\bullet$ & $\bullet$ & validator + memory (no refine loop)                    \\
    R+X & $\bullet$ &         & $\bullet$ &         & retrieval + validator (no refine loop)                 \\
    S+X &         & $\bullet$ & $\bullet$ &         & refine loop (with xmllint) + validator                 \\
    R+S+X+M & $\bullet$ & $\bullet$ & $\bullet$ & $\bullet$ & full stack                                     \\
    \midrule
    S+X+M    &         & $\bullet$ & $\bullet$ & $\bullet$ & factorial completion (predicted main-effects best) \\
    SE-prose &         & $\bullet$ & $\bullet$ & $\bullet$ & S+X+M settings, v3 prose + v3 cheatsheet for m1u   \\
    SE       &         & $\bullet$ & $\bullet$ & $\bullet$ & self-evolved v3 plugin (\S\ref{subsec:overview})   \\
    \bottomrule
  \end{tabular}
\end{table}

\section{Bottleneck-analysis pipeline (full)}
\label{app:bottleneck}

\textbf{Stage 1 (LLM-free extractor).} For each (cell, seed, task), the extractor walks the recursive \texttt{treesim\_detail} tree from the scorer and ranks subtrees by impact $= (1 - \mathrm{score}) \cdot (n_{\mathrm{gt\_children}} + 1)$, returning the top-$K$. It also extracts trajectory features from \texttt{events.jsonl}: tool counts, file re-reads, xmllint MCP calls, edit churn, grep/glob query distributions.

\textbf{Stage 2 (\texttt{deepseek-v4-flash} per-task classifier).} The extractor's output, plus a focused GT-vs-generated XML excerpt for the worst-scoring section when its score is below 0.7, is sent to \texttt{deepseek-v4-flash} with a strict-JSON output schema. The model returns: \texttt{failure\_category} $\in \{$\texttt{missing\_block}, \texttt{extra\_block}, \texttt{hallucinated\_extras}, \texttt{structural\_mismatch}, \texttt{bad\_attribute\_value}, \texttt{partial\_implementation}, \texttt{wrong\_constitutive}, \texttt{no\_failure}$\}$, a \texttt{primary\_failure\_section}, a one-sentence \texttt{root\_cause}, a \texttt{trajectory\_evidence} citation, and a \texttt{would\_have\_helped} hypothesis.

\textbf{Stage 3 (\texttt{deepseek-v4-pro} synthesis).} Per-cell category and section distributions, section ``failure weight'' $= \sum_{\mathrm{tasks}}(1 - \mathrm{TreeSim})$ grouped by primary failing section, and per-task baseline-vs-best deltas are sent to \texttt{deepseek-v4-pro} with an instruction to write a paper-grade synthesis anchored on a chosen baseline/best pair. Total cost across all evaluations was ${\sim}650$ flash calls plus 4 pro narratives, ${\sim}\$5\text{--}7$ in DeepSeek API spend.

Code: \texttt{scripts/bottleneck/\{extract,llm\_per\_task,aggregate\}.py}.

\subsection{Bottleneck failure-category counts (full)}
\label{app:bottleneck-counts}

\begin{table}[h]
  \caption{Failure-category counts per cell. val panel uses $n=51$ tasks per cell ($3{\times}17$); held-out-eval panel uses $n=29$--$30$ ($3{\times}10$ minus a small number of LLM-judge parse-failures).}
  \label{tab:bottleneck}
  \centering
  \footnotesize
  \resizebox{0.98\textwidth}{!}{%
  \begin{tabular}{l|ccccc|cccccc}
    \toprule
    & \multicolumn{5}{c|}{\textbf{val ($n=51$)}} & \multicolumn{6}{c}{\textbf{held-out-eval ($n=30$)}} \\
    \textbf{Category} & Vanilla & S+M & X+M & S+X & SE & Vanilla & X+M & S+X & S+X+M & SE-prose & SE \\
    \midrule
    bad\_attribute\_value     & 12 & 12 & 11 & 15 & 9  & 5 & 7 & 5 & 8 & 0 & 0 \\
    extra\_block              &  9 &  8 & 11 &  5 & 10 & 9 & 4 & 4 & 6 & 7 & 5 \\
    hallucinated\_extras      &  4 &  8 &  7 &  - &  - & 0 & 0 & 0 & 0 & 3 & 4 \\
    missing\_block            &  6 &  4 &  3 &  2 &  - & 4 & 5 & 7 & 7 & 7 & 4 \\
    partial\_implementation   &  7 &  9 &  6 & 10 &  9 & 1 & 0 & 4 & 3 & 3 & 5 \\
    structural\_mismatch      &  6 &  8 &  7 & 11 &  - & 8 & 6 & 6 & 5 & 3 & 5 \\
    \bottomrule
  \end{tabular}}
\end{table}

\section{Per-task held-out-eval results}
\label{app:per-task}

\begin{table}[h]
  \caption{Per-task held-out-eval TreeSim (mean across 3 seeds), sorted ascending by Vanilla score. The aggregate $+0.069$ Vanilla$\to$SE gain is concentrated in the top two non-universal-failure tasks (rows 2--3); the remaining seven tasks show within-noise differences across cells.}
  \label{tab:per-task-icl10}
  \centering
  \footnotesize
  \resizebox{0.98\textwidth}{!}{%
  \begin{tabular}{lrrrrrr}
    \toprule
    \textbf{Task} & \textbf{Vanilla} & \textbf{X+M} & \textbf{S+X} & \textbf{S+X+M} & \textbf{SE-prose} & \textbf{SE} \\
    \midrule
    \texttt{TutorialHydraulicFractureWithAdvancedXML}      & 0.013 & 0.013 & 0.013 & 0.013 & 0.013 & 0.013 \\
    \texttt{AdvancedExampleThermoPoroElasticWellbore}      & 0.355 & 0.680 & 0.681 & 0.708 & 0.743 & \textbf{0.761} \\
    \texttt{ExampleProppantTest}                           & 0.541 & 0.810 & 0.809 & 0.799 & 0.817 & \textbf{0.825} \\
    \texttt{ExampleIsothermalHystInjection}                & 0.755 & 0.747 & 0.751 & 0.750 & 0.769 & 0.717 \\
    \texttt{AdvancedExampleCasedThermoElasticWellbore}     & 0.847 & 0.807 & 0.923 & 0.919 & 0.877 & 0.886 \\
    \texttt{ExamplesingleFracCompression}                  & 0.891 & 0.887 & 0.904 & 0.931 & 0.929 & 0.928 \\
    \texttt{ExampleVerticalPoroElastoPlasticWellbore}      & 0.909 & 0.903 & 0.906 & 0.891 & 0.834 & 0.944 \\
    \texttt{ExampleMCCWellbore}                            & 0.935 & 0.924 & 0.908 & 0.905 & 0.905 & 0.941 \\
    \texttt{AdvancedExamplePureThermalDiffusionWellbore}   & 0.963 & 0.922 & 0.956 & 0.947 & 0.864 & 0.880 \\
    \texttt{AdvancedExampleViscoExtendedDruckerPrager}     & 0.986 & 0.991 & 0.963 & 0.964 & 0.999 & 0.996 \\
    \bottomrule
  \end{tabular}}
\end{table}

\section{Resolution-IV main effects on val}
\label{app:main-effects}

The eight factorial cells on val yield the per-factor main effects below, estimated as the mean TreeSim difference between the four cells with the factor on and the four with it off. R is the only factor whose magnitude exceeds the seed-noise floor; X and M are positive but small; S is essentially zero. On val alone the design implication is ``don't add RAG; the rest doesn't matter.'' The held-out-eval results in \S\ref{sec:results} show this conclusion is benchmark-specific.

\begin{table}[h]
  \caption{Resolution-IV main effects on val TreeSim (\texttt{deepseek-v4-flash}, $n=3$ per cell).}
  \label{tab:main-effects}
  \centering
  \small
  \begin{tabular}{lr}
    \toprule
    \textbf{Factor} & \textbf{$\Delta$ TreeSim} \\
    \midrule
    R (RAG)                & $-0.032$ \\
    S (Stop-hook)          & $-0.003$ \\
    X (xmllint MCP)        & $+0.007$ \\
    M (memory cheatsheet)  & $+0.004$ \\
    \bottomrule
  \end{tabular}
\end{table}

\section{Cross-model and cross-harness: full panels}
\label{app:cross-model-detail}
\label{app:cross-model}

\textbf{Setup.} Cross-model: $n=1$ seed of Vanilla, X+M, SE on val with two backbones (\texttt{minimax-m2.7}, \texttt{google/gemini-3-flash-preview}) via OpenRouter, harness held at Claude Code, default-off harness env (\texttt{GEOS\_HOOK\_POSTTOOLUSE} unset, autocamp-experiment-state parity). Cross-harness: $n=3$ seeds of Vanilla and X+M on val with \texttt{deepseek-v4-flash}, harness varied between Claude Code and OpenHands. Headline DSv4 numbers from Table~\ref{tab:main-results} included for comparison.

\begin{table}[h]
  \caption{Full cross-model and cross-harness panel. Counts in parentheses are task failures out of 17.}
  \label{tab:cross-cutting-full}
  \centering
  \small
  \resizebox{0.98\textwidth}{!}{%
  \begin{tabular}{llccc}
    \toprule
    \textbf{Harness} & \textbf{Backbone} & \textbf{Vanilla} & \textbf{X+M} & \textbf{SE} \\
    \midrule
    CC & \texttt{deepseek-v4-flash} ($n=3$) & $0.910 \pm 0.024$ & $0.921 \pm 0.007$ & $0.919 \pm 0.020$ \\
    CC & \texttt{minimax-m2.7} ($n=1$)        & $0.821$ ($1$ fail) & $0.867$ ($0$ fails) & $0.861$ ($0$ fails) \\
    CC & \texttt{gemini-3-flash-preview} ($n=1$) & $0.768$ ($0$ fails) & $0.797$ ($0$ fails) & $0.757$ ($1$ fail) \\
    OH & \texttt{deepseek-v4-flash} ($n=3$) & $0.856 \pm 0.061$ & $0.881 \pm 0.023$ & -- \\
    \bottomrule
  \end{tabular}}
\end{table}

\textbf{R-factor cross-model corroboration} (older runs). On a 35-task superset that overlaps val, the retrieval plugin lifts $+0.175$ on \texttt{deepseek-v3.2} (paired 35 tasks) and $+0.115$ on \texttt{minimax-m2.7} (paired 15 tasks). These older runs predate the bug discovery; the headline result above (X+M and SE × val) supersedes them.

\textbf{Harness-less floor.} A one-shot direct-prompting baseline on \texttt{minimax-m2.7} (\texttt{scripts/harnessless\_eval.py}; \texttt{deepseek-v4-flash} not yet run in this configuration) reaches TreeSim $= 0.333$ on val. We treat this as a model-class lower bound rather than a per-model floor.

\textbf{Memory-as-retrieval negative result.} An earlier version of our system exposed primer-like content via an MCP \texttt{memory\_lookup} tool with embedding top-$k$ retrieval over the distilled corpus. Across every test-set run in which this tool was available (the A4, A4$'$, A5, M3-g conditions of the prior campaign), the agent called the tool zero times. The tool was verified functional. Delivering the same content through the always-on \texttt{--append-system-prompt} channel (i.e., the M factor in this paper) is what produced any lift. We retain this as a clean negative result on retrieval-based memory for this task--model class.

\section{OpenFOAM transfer study}
\label{app:openfoam-transfer}

\paragraph{Motivation.}
The main paper studies SIGA on GEOS, where the executable interface is a structured XML deck with an explicit XSD schema. To test whether the same adapter recipe transfers beyond GEOS and beyond XML-backed simulator interfaces, we ran a small OpenFOAM case-authoring companion study in a separate repository path (\texttt{repo3\_openfoam}). We treat the result as transfer evidence rather than a second benchmark because the setup differs in both metric and scale.

\paragraph{OpenFOAM adaptation of \texorpdfstring{$R/S/X/M$}{R/S/X/M}.}
We ported the same four binary factors to OpenFOAM case authoring. \textbf{R} replaced the GEOS RAG collections with three ChromaDB collections over OpenFOAM tutorial structure, detailed case snippets, and command/help text. \textbf{M} replaced the GEOS memory cheatsheet with a lightweight always-on primer describing the standard OpenFOAM case skeleton (\texttt{0/}, \texttt{constant/}, \texttt{system/}) and the relevant tutorial/source-tree locations. \textbf{S} replaced the GEOS XML stop-hook with an OpenFOAM stop-hook that blocks turn completion if required files are missing or if generated dictionaries fail heuristic structural checks. \textbf{X} replaced the \texttt{xmllint} MCP with an agent-callable OpenFOAM validator using the same case checks as the hook. Because OpenFOAM lacks a canonical XSD-style schema for the benchmark dictionaries, the OpenFOAM validator checks required-file presence, balanced delimiters, \texttt{FoamFile} headers, and key dictionary sections rather than schema compliance.

\paragraph{Benchmark and metric.}
The OpenFOAM subset contains 5 tasks sampled from a FoamGPT-derived benchmark: \texttt{boundaryWallFunctionsProfile}, \texttt{Grossetete}, \texttt{helmholtzResonance}, \texttt{externalCoupledCavity}, and \texttt{damBreakWithObstacle}. Each task specifies a set of required files. We score a run using a file-text-and-coverage metric: each missing required file receives score 0; present files are compared to ground truth via normalized text similarity; and the case score is
\[
\mathrm{Score} = 0.7 \cdot \mathrm{mean\_similarity} + 0.3 \cdot \mathrm{coverage}.
\]
This is analogous to failures-as-zero reporting but is not TreeSim.

\paragraph{Foam-Agent baseline.}
The OpenFOAM-native baseline is Foam-Agent~\citep{foamagent2025}. We emphasize a caveat that matters for interpretation: Foam-Agent's native workflow is not lint-only, but a fuller plan--write--execute--review loop. The stable comparison available in our environment used \texttt{execution\_mode=lint\_only}; execute-mode runs failed to yield usable benchmark outputs. We therefore compare against \emph{Foam-Agent} and treat the result as a constrained baseline rather than a full head-to-head against Foam-Agent's intended execution-coupled workflow.

\begin{table}[t]
  \caption{OpenFOAM transfer-study summary on the 5-task subset. Mean score is the file-text-and-coverage metric described above. \emph{Full coverage} is the number of tasks for which all required files were produced.}
  \label{tab:openfoam-summary}
  \centering
  \small
  \resizebox{0.98\textwidth}{!}{%
  \begin{tabular}{lccccc}
    \toprule
    \textbf{Cell} & \textbf{Mean score} & \textbf{$\Delta$ vs Vanilla} & \textbf{$\Delta$ vs Foam-Agent$_{\text{lint}}$} & \textbf{Full coverage} & \textbf{Wall s} \\
    \midrule
    Vanilla                  & 0.466 & --    & $-0.103$ & 3/5 & 1921 \\
    R$+$M                    & 0.736 & $+0.270$ & $+0.167$ & 5/5 & 1276 \\
    S$+$M                    & 0.787 & $+0.321$ & $+0.218$ & 5/5 & 1482 \\
    R$+$S                    & \textbf{0.871} & \textbf{$+0.405$} & \textbf{$+0.302$} & \textbf{5/5} & \textbf{380} \\
    X$+$M                    & 0.712 & $+0.246$ & $+0.143$ & 5/5 & 1644 \\
    R$+$X                    & 0.145 & $-0.321$ & $-0.424$ & 1/5 & 778 \\
    S$+$X                    & 0.849 & $+0.383$ & $+0.280$ & \textbf{5/5} & 1020 \\
    R$+$S$+$X$+$M            & 0.862 & $+0.396$ & $+0.294$ & \textbf{5/5} & 1419 \\
    S$+$X$+$M                & 0.822 & $+0.356$ & $+0.253$ & \textbf{5/5} & 1147 \\
    Foam-Agent$_{\text{lint}}$ & 0.569 & $+0.103$ & --       & 3/5 & 547 \\
    \bottomrule
  \end{tabular}}
\end{table}

\paragraph{Headline.}
The main OpenFOAM result is qualitative rather than statistical: the same component that matters most in GEOS reliability is again the strongest signal here. The best OpenFOAM cell is R+S at 0.871, well above vanilla Claude Code (0.466) and Foam-Agent$_{\text{lint}}$ (0.569). More importantly, every $\mathrm{S}$-enabled cell achieves full required-file coverage and no zero-score tasks, while vanilla covers only 3/5 tasks and R+X only 1/5.

\paragraph{Factor-style readout.}
Using the 8-cell OpenFOAM subset only and reporting single-run descriptive effects, the factor-style main-effect analogue is:
\[
R: -0.050,\quad
S: +0.328,\quad
X: -0.073,\quad
M: +0.192
\]
on mean score. This again points to $\mathrm{S}$ as the dominant reliability intervention and suggests that optional retrieval or optional validation, without a hard end-of-turn gate, are weaker.

\begin{table}[t]
  \caption{Per-task OpenFOAM transfer-study scores. The catastrophic failures are concentrated in non-$\mathrm{S}$ cells, especially Vanilla and R+X.}
  \label{tab:openfoam-per-task}
  \centering
  \footnotesize
  \resizebox{0.98\textwidth}{!}{%
  \begin{tabular}{lccccc}
    \toprule
    \textbf{Cell} & \textbf{boundaryWallFunctionsProfile} & \textbf{Grossetete} & \textbf{helmholtzResonance} & \textbf{externalCoupledCavity} & \textbf{damBreakWithObstacle} \\
    \midrule
    Vanilla                  & 0.751 & 0.817 & 0.000 & 0.762 & 0.000 \\
    R$+$M                    & 0.650 & 0.817 & 0.681 & 1.000 & 0.532 \\
    S$+$M                    & 0.860 & 0.966 & 0.788 & 0.595 & 0.723 \\
    R$+$S                    & 0.907 & 0.968 & 0.858 & 0.900 & 0.723 \\
    X$+$M                    & 0.719 & 0.820 & 0.727 & 0.762 & 0.531 \\
    R$+$X                    & 0.724 & 0.000 & 0.000 & 0.000 & 0.000 \\
    S$+$X                    & 0.965 & 0.870 & 0.787 & 0.899 & 0.723 \\
    R$+$S$+$X$+$M            & 0.964 & 0.967 & 0.760 & 0.899 & 0.722 \\
    S$+$X$+$M                & 0.807 & 0.788 & 0.850 & 0.943 & 0.722 \\
    Foam-Agent$_{\text{lint}}$ & 0.657 & 0.165 & 0.649 & 0.636 & 0.736 \\
    \bottomrule
  \end{tabular}}
\end{table}

\paragraph{Zero-score failures.}
The zero scores in the OpenFOAM study are failures-as-zero in the literal sense: required output files were missing. Vanilla drops to 0 on \texttt{helmholtzResonance} and \texttt{damBreakWithObstacle}; R+X drops to 0 on 4 of 5 tasks. No $\mathrm{S}$-enabled cell exhibits a zero-score failure. This is the clearest OpenFOAM analogue to the GEOS reliability story: the stop hook prevents the agent from ending the turn with structurally incomplete deliverables.

\paragraph{Interpretation.}
We do not claim the OpenFOAM result has the same evidentiary weight as the GEOS benchmark. It is smaller, single-seed, and scored differently. What it does show is that the SIGA recipe is not obviously tied to XML or to GEOS-specific assets. The transferable piece is the workflow logic: retrieval and procedural guidance matter somewhat, but the dominant and most portable intervention is forced end-of-turn verification that prevents silent incompleteness.

\section{Future work: open follow-ups}
\label{app:future-work}

\textbf{Multi-seed cross-model.} Repeat the cross-model panel (Vanilla, X+M, SE on minimax and gemini) with $n=3$ seeds to tighten effect-size estimates. Cost is the limiting factor: at OpenRouter list pricing, gemini-3-flash-preview is the dominant share (${\sim}3\times$ DSv4 input, ${\sim}11\times$ output per million tokens). Single-seed already establishes that the X+M ranking transfers; multi-seed sharpens the magnitude.

\textbf{Multi-seed Foam-Agent in execute mode.} The OpenFOAM transfer comparison currently uses Foam-Agent in lint-only mode because execute-mode runs failed to yield usable benchmark outputs in our environment. Re-running Foam-Agent in its native execute-and-review mode (with $n=3$ seeds on the same 5-task subset) would convert the constrained baseline into a definitive head-to-head. We list this as the natural follow-up to the cross-simulator transfer evidence.

\textbf{Larger OpenFOAM benchmark.} Scale the OpenFOAM transfer subset from 5 to 20+ tasks, drawn across the full FoamGPT distribution, and add multi-seed runs. This converts the qualitative ``transfer evidence'' framing into a second benchmark suitable for quantitative effect-size estimation.

\textbf{OpenCode harness integration.} Build a thin wrapper for OpenCode (no extant integration) and re-run Vanilla and X+M-equivalent. We outlined the implementation cost as $\sim$4--8h.

\textbf{Execution-correctness ladder.} Run a sample of agent-produced decks through actual GEOS execution to convert TreeSim into a runnability metric. Even a small panel (5 tasks × 2 cells × 1 seed) would convert the structural-similarity metric into a ladder.

\textbf{Cross-section consistency hook} (from the bottleneck analysis). Implement an adapter that validates \texttt{<ElementRegion materialList>} entries against \texttt{<Constitutive>} block names, then re-run as a new cell. This converts the bottleneck-analysis-as-observation into adapter-design-as-output.

\section{Implementation details}
\label{app:impl}

\subsection{Hook wiring}
Historical note: prior to 2026-04-21T12:08Z, \texttt{hooks.json} had the wrong schema and \texttt{run\_experiment.py} never passed \texttt{--plugin-dir}, so early ``hook-on'' runs did not actually load the hook. All post-fix results in this paper use a verified-wired hook.

\subsection{Primer artefacts and hygiene}
The memory cheatsheet is at \texttt{plugin/memory\_primer\_m1u.md}\footnote{filename pending final confirmation}. The hygiene audit is at \texttt{scripts/memory/hygiene\_audit.py}; the API-contamination check is at \texttt{scripts/memory/check\_api\_contamination.py}.

\subsection{Harness-less baseline}
\label{subsec:harnessless}
Harness-less eval: \texttt{scripts/harnessless\_eval.py}. Inputs: one held-out ICL demonstration (task plus ground-truth XML), one call per task, with an inline-XML protocol replacing all file-system instructions. System prompt = \texttt{run/AGENTS.md} with file-system paragraphs stripped and replaced by the inline-XML protocol. Temperature 0.2, \texttt{max\_tokens} $=$ 16384, 600-s per-request timeout, 8-worker thread pool. Result on val (\texttt{minimax-m2.7}): TreeSim $= 0.333$ (16 parsed, 1 silent provider drop counted as zero). The vanilla Claude Code harness on the same model recovers ${+}0.164$ over this floor; any cell on \texttt{deepseek-v4-flash} reaches a different absolute level ($\geq 0.857$). In both cases the harness contributes substantially even before any GEOS-specific adaptation.

\section{Agent autonomy companion: protocol}
\label{app:autonomy}

This appendix complements \S\ref{subsec:autonomy} with the protocol and intermediate diagnostics for the autonomy companion study.

\paragraph{Difficulty tiers.} We tag each ground-truth parameter into one of four tiers, following a domain-expert-reviewed taxonomy: T1 \emph{software defaults} (output format, restart frequency, log levels, boolean flags), T2 \emph{standard numerics} (Newton tolerance, linear solver, time-step limits, discretisation choice, element type), T3 \emph{domain-inferable} (densities, viscosities, porosities, permeabilities, Biot coefficients, standard relperm exponents), and T4 \emph{problem-defining} (geometry, well locations, applied loads, simulation duration, prescribed history tables). Medium difficulty omits T1+T2; Hard additionally omits T3. T4 is preserved verbatim at every level.

\paragraph{Spec generation.} Each task's \texttt{instructions.txt} is rewritten by \texttt{deepseek-v4-pro} (\texttt{scripts/relax\_specs.py}) given the original brief and the tier definitions, in two passes (Medium, Hard). The model is required to drop only, never to invent or alter values, and to keep external-table references (e.g., \texttt{*.geos} input files for boundary-history tables) at every level. Each rewrite is paired with a JSON record listing the values dropped, the tier of each, and a hygiene check that canonicalises numeric tokens (handling LaTeX \texttt{\$1.0\textbackslash times10\^{}\{-4\}\$}, scientific notation, unicode superscripts) and verifies that no canonicalised dropped value appears in the rewrite text except where the same canonical also appears in a kept-T4 value (``shared''). Across all 16 (task $\times$ level) pairs, hygiene flags zero leaks after the LaTeX-aware pass and 1 shared-only ambiguity (Mandel/Medium, where the dropped solver tolerance \texttt{1e-4} canonicalises to the same value as the first entry of the kept T4 prescribed-displacement time table). The 16 rewrites are committed and frozen before any run starts.

\paragraph{Drop volume (8 tasks total).}

\begin{table}[h]
  \caption{Volume of specification dropped at each difficulty level. T1=software defaults, T2=standard numerics, T3=domain-inferable physical values, T4=problem-defining (kept verbatim at every level). Char count is summed across all 8 tasks.}
  \label{tab:autonomy-drop}
  \centering
  \small
  \begin{tabular}{lcccc}
    \toprule
    \textbf{Level} & \textbf{Char drop} & \textbf{Values dropped} & \textbf{T1 / T2 / T3 split} \\
    \midrule
    Medium & $-20.5\%$ & 89  & 27 / 62 / 0 \\
    Hard   & $-50.5\%$ & 184 & 26 / 53 / 105 \\
    \bottomrule
  \end{tabular}
\end{table}

\paragraph{Supervisor channel.} The interactive cells expose an MCP server (\texttt{plugin/scripts/supervisor\_mcp.py}) with a single tool, \texttt{consult\_supervisor(question: str)}. The handler invokes \texttt{deepseek-v4-flash} with a system prompt that injects the \emph{full} original \texttt{instructions.txt} for the current task and the rule ``answer concisely using only information present in the specification; if the answer is not in the specification, say so plainly; do not invent; do not volunteer information the agent did not ask about; speak in the same scientific language the specification uses (no GEOS XML tag names unless the agent named them first)''. The full original brief is mounted at a fixed container path outside \texttt{/workspace}; the path is passed via the per-MCP-server env block in the explicit \texttt{--mcp-config}, not via \texttt{docker -e}. Every consultation appends a structured record (question, answer, latencies, token counts) to \texttt{/workspace/supervisor\_calls.jsonl} for audit. We programmatically scanned all 32 \texttt{events.jsonl} streams for any \texttt{Read} or \texttt{Bash} tool input referring to the supervisor's spec path; zero hits. The path leaks via the \texttt{supervisor\_stats} introspection tool's return value (a known minor leak; should be redacted in any future iteration).

\paragraph{Prompt variants.} Two are reported. \textbf{V0} (the default) describes the channel in both the \texttt{consult\_supervisor} docstring and a system-prompt addendum as something to use ``when a value or design decision you need is missing from the brief AND you cannot reasonably infer it from GEOS conventions, GEOS example simulations, or standard geophysics practice. Each call costs the researcher's time, so prefer to infer when you can.'' \textbf{V1 (neutral)} drops the prefer-to-infer language and treats infer-vs-ask as peer paths: ``values that are missing can be inferred from GEOS conventions and analogous examples, OR you may ask the researcher; choose whichever path is more reliable.'' V0 was selected by 1/32 trials; V1 was selected by 1/32 trials.

\paragraph{Diagnostic on Mandel/Hard: on-disk findability.} For each T2/T3 value dropped from the Mandel/Hard rewrite (26 values), we test whether the same canonicalised numeric token appears anywhere in the GEOS-shipped example XMLs that the agent is allowed to read (excluding the GT files for the current task, which are blocked by contamination): 15 of 26 are findable (e.g., \texttt{0.001 Pa\,s} fluid viscosity in 210 other examples; \texttt{1e-12 m\textsuperscript{2}} permeability in 51; reference porosity \texttt{0.375} in 8). One value (bulk modulus 66.667 MPa, an unusually rounded value) has zero hits; the rest are non-numeric T2 directives (subcycling, line-search action, discretisation choice) the agent can plausibly transfer from analogous benchmarks. We interpret this as a structural reason for the consultation rate observed in \S\ref{subsec:autonomy}.

\paragraph{Where the agent looks.} Aggregate \texttt{Read}/\texttt{Glob} call counts into \texttt{/geos\_lib/inputFiles/} across all 8 tasks per cell range from $142$ (X+M, Medium) to $404$ (Vanilla, Medium) per cell. The agent is not idle; it is grepping the example library.

\paragraph{Cost and wall.} 64 main runs at \texttt{deepseek-v4-flash} cost ${\sim}\$4.20$ DeepSeek API spend at full off-peak pricing (input + cache-read + output, summed from \texttt{events.jsonl}) and ran in $\sim$2h wall at \texttt{workers=4}. 16 \texttt{deepseek-v4-pro} spec-rewrite calls cost ${\sim}\$0.50$. The 2 supervisor consultations cost ${\sim}\$0.001$. The 32-run V1 rerun added ${\sim}\$2$ and ${\sim}55$min wall.

\paragraph{Immediate follow-ups.} (a) A second seed per cell to harden the $n=1$ point estimates ($\sim\$3$, $\sim$2h). (b) A no-confound F0 control (supervisor MCP wired without the plugin loader) so the F0 vs F0+supervisor comparison is not contaminated by tool-list-shape effects. (c) An aggressive contamination block at the physics-family level (e.g., for \texttt{ExampleMandel}, block all \texttt{*Mandel*} and \texttt{*Poroelastic*} input files, not just the GT basenames). The Mandel/Hard diagnostic predicts (c) is what would actually move the consultation rate; it directly tests whether the on-disk example library is the binding constraint.

\section{Human baseline: protocol and browser-history breakdown}
\label{app:human-browser}

This appendix complements \S\ref{subsec:human-baseline} with the protocol and the per-domain navigation breakdown.

\paragraph{Protocol.} Two geoscience-domain-expert volunteers (Expert 1, Expert 2; grad-level geoscientists, new to GEOS-the-software) attempted the \texttt{buckleyLeverettProblem} task (1D Buckley--Leverett CO$_2$/brine displacement; vanilla Claude Code on \texttt{deepseek-v4-flash} reaches deck-level TreeSim $0.751 \pm 0.016$ on this task at $n=3$ seeds) under a one-hour timeslot. Participants received the same system primer the agent receives (\texttt{run/AGENTS.md} including the GEOS Primer), a working directory mirroring the agent's container layout (empty \texttt{inputs/}, empty \texttt{outputs/}), and read access to a filtered copy of the GEOS source tree mounted at the same path the agent sees. The blocked file list was identical to the agent's contamination block (the three ground-truth XMLs and the tutorial \texttt{Example.rst}). Participants were instructed not to use LLM chatbots, code-completion services, or general internet search; they were free to consult the GEOS Sphinx documentation and the GEOS GitHub repository, and to use \texttt{grep}/\texttt{find} on their own machine. Wall-clock was measured from spec-opened to deck-final; participants self-reported. We additionally collected the browser history for each participant's session (Expert 1: a CSV export from a Chrome history extension; Expert 2: a JSON export from a Firefox history extension) and filtered to the day of the session. Expert 1 subsequently returned to the assignment with no time cap and produced both required files; we report the catch-up session below.

\paragraph{Per-domain breakdown.}

\begin{table}[h]
  \caption{Browser-navigation breakdown during the one-hour authoring session, per participant. ``GEOS docs'' = visits to the GEOS Sphinx documentation host (\texttt{geosx-geosx.readthedocs-hosted.com}) including all subpages. ``GEOS GitHub'' = visits to \texttt{github.com/GEOS-DEV/GEOS}. ``Other'' = unrelated tabs that remained open during the session. No participant visited any LLM chatbot, Stack Overflow, or general scientific-paper site.}
  \label{tab:human-browser}
  \centering
  \small
  \resizebox{0.98\textwidth}{!}{%
  \begin{tabular}{lcccccc}
    \toprule
    \textbf{Participant} & \textbf{Wall (min)} & \textbf{Total visits} & \textbf{GEOS docs} & \textbf{GEOS GitHub} & \textbf{Search} & \textbf{Other (Slack, etc.)} \\
    \midrule
    Expert 1 (1\,h)                & $48.2$       & $29$  & $20$ & $5$  & $3$ & $1$ \\
    Expert 1 (extended, +catch-up) & ${\sim}180$  & $106$ & $89$ & $21$ & $6$ & $7$ \\
    Expert 2 (1\,h)                & $46.7$       & $73$  & $54$ & $11$ & $5$ & $3$ \\
    \bottomrule
  \end{tabular}}
\end{table}

\paragraph{Top GEOS pages visited.} Across both participants, the most-visited GEOS pages were the \texttt{CompositionalMultiphaseFlow} solver reference, the \texttt{EventManager} page (which both participants flagged as the source of difficulty in their post-session notes), the \texttt{Outputs} schema, and the GEOS \texttt{Tutorials/Index} landing page. Both participants browsed the \texttt{compositionalMultiphaseFlow/dbc/buckleyLeverett\_1d/} sibling deck on GEOS GitHub, a separate 1D Buckley--Leverett deck (DBC variant) that is not the ground truth and is not blocked. In post-session notes, Expert 2 reported that \texttt{Outputs} and \texttt{Events} setup ``took forever to try and understand because it is not very clear and pretty unintuitive''; Expert 1 reported using the DBC deck as a structural template, then adjusting mesh/run-time/solver parameters to match the spec. Within the 1h budget, neither attempted to write the second required file (\texttt{buckleyLeverett\_benchmark.xml}); Expert 1 produced it during the extended catch-up session.

\paragraph{Catch-up session: what changed.} Expert 1's catch-up session adds $77$ navigations on top of the original $29$, of which $69$ are GEOS-internal (Sphinx + doxygen + GitHub). The qualitative strategy is unchanged (Sphinx-prose-driven authoring with sibling decks pulled from GitHub for structural templates), but the docs surface broadens: heavy re-reading of \texttt{EventManager}, \texttt{Outputs}, and \texttt{TasksManager} (consistent with Expert 1's written note that ``2/3 of this time'' went to the outputs/events portion of the deck); three doxygen visits to \texttt{PhaseVolumeFractionKernel.hpp} (Expert 1 reports having to ``look at the source code to get specific inputs for the prompt, specifically the \texttt{fieldName} for the .hdf5 output''); and several visits to an unrelated \texttt{compositionalMultiphaseWell/benchmarks/Egg} deck used as a structural template for outputs/events/tasks blocks. The agent's behaviour on this same task does not change between the two ``human sessions''; it does not visit any of these pages, and it solves the outputs/events portion in the same ${\approx}\,7\,$min envelope.

\paragraph{Disclosure: one ChatGPT navigation in the catch-up.} Expert 1's catch-up history shows one navigation to \texttt{chatgpt.com} on 05/06/26 at 11:21 (conversation title ``Linux file search tips''). Adjacent entries (Google searches for \texttt{libmainInterface.so} shared-library errors, visits to UCI's RCIC cluster documentation) place this visit during a post-authoring attempt to run the produced deck on the cluster, ${\sim}90\,$min after the bulk of the XML editing. We disclose it as a deviation from our ``no LLM chatbots'' instruction; we judge that it did not affect the deck content and have left the extended-run scores in Table~\ref{tab:human-baseline}. If anything this deviation strengthens the human-baseline anchor: an LLM-augmented domain-expert with no time cap reaches deck-level TreeSim $0.931$, a number very close to the SIGA agent's ceiling on this task. Expert 2's session contained no such navigations.

\paragraph{Agent-side counterpart.} Cross-run file-access analysis of $7$ vanilla-Claude-Code runs on the same task (\texttt{scripts/analysis/analyze\_file\_access.py}) gives means of ${\approx}\,14$ unique files per run, ${\approx}\,11$ XML reads from \texttt{/geos\_lib/inputFiles/}, ${\approx}\,7$ \texttt{Grep} calls, ${\approx}\,5$ \texttt{Glob} calls; plugin variants drop unique-file reads to ${\approx}\,4$ and replace \texttt{Grep}/\texttt{Glob} with ${\approx}\,2$ structured retrieval calls. The cumulative file/document surface the humans access (Sphinx pages plus GitHub-rendered XMLs) is roughly comparable in count to the agent's file-read surface, but they are different files: the humans read prose narrating the simulator's vocabulary, the agent reads concrete XMLs that exemplify it.

\paragraph{Outputs.} Scoring scripts: \texttt{scripts/score\_human\_baseline.py}, \texttt{scripts/analyze\_human\_browser\_history.py}. Browser-history exports and submitted XMLs at \texttt{data/human\_baseline/}. A narrative summary is at \texttt{docs/2026-05-04\_human-baseline-browser-analysis.md}.

\paragraph{Caveats.} $n=2$ on a single task in a one-hour budget is a calibration anchor, not a study of human authoring competence. The participants are domain-expert subsurface modellers, not GEOS power users; a longer-time budget or a participant pool of long-time GEOS users would likely change the absolute level. The extended-budget sanity check on Expert 1 raises the deck-level number to $0.931$ at $\sim\!3\,$h wall-clock, evidence that the 1-hour budget, not domain-expert capacity, is the binding constraint on the deck-level shortfall. The GEOS-expert estimate (\S\ref{subsec:human-baseline}) brackets ``simple Buckley--Leverett'' at ${<}30\,$min for a power user starting from a known-good deck, locating the agent (${\approx}\,7\,$min) inside that bracket. We treat the (Expert 1 at 1h, Expert 2 at 1h) pair as the headline budget-matched comparison; the (Expert 1 extended) pair as orthogonal anchors that contextualise it. We treat the overall result as an existence-of-effect on three points: (i) the absolute TreeSim level on \texttt{buckleyLeverettProblem} is roughly $0.78$--$0.81$ at the file level under a one-hour domain-expert budget; (ii) the file-usage strategy a human reaches for under documentation pressure is qualitatively different from the agent's source-tree-example strategy; (iii) both participants ran out of time on a single deck before producing the second of the two required files within the 1h budget, while the agent reaches a comparable file-level number on the full two-file task in roughly seven minutes.

\section{Example GEOS deck}
\label{app:geos-example}

We reproduce an abridged form of the \texttt{buckleyLeverettProblem} ground-truth deck (the same task used in the human-baseline study, \S\ref{subsec:human-baseline}) to make the cross-section constraints described in \S\ref{sec:background} concrete. The deck is split across two XML files in the GEOS convention: \texttt{buckleyLeverett\_base.xml} declares physics, materials, and boundary conditions; \texttt{buckleyLeverett\_benchmark.xml} \texttt{<Included>}s the base file and adds the mesh, geometry, and event timeline. Comments and verbose attribute defaults are elided for space; the original is 172 + 61 lines.

\paragraph{base file: physics, materials, boundary conditions.}
\promptinput{assets/buckleyLeverett_base.xml}

\paragraph{benchmark file: mesh, geometry, events.}
\promptinput{assets/buckleyLeverett_benchmark.xml}

\paragraph{Cross-section constraints visible in this deck.}
The canonical \emph{ten}-section GEOS schema collapses here to a smaller set of blocks (no \texttt{Functions}, no separate \texttt{Mesh} in the base file) because Buckley--Leverett is at the easy end of the bench; harder tasks add all ten. Even on this deck, deck authoring requires several program-level constraints to hold simultaneously, and these are exactly the constraints the bottleneck analysis (\S\ref{subsec:bottleneck-results}) flags as adapter-resistant when violated:

\begin{itemize}
\item \textbf{Region-name matching across sections.} The string \texttt{region} is declared once in \texttt{<ElementRegions>} and is re-used as a member of \texttt{targetRegions} on the \texttt{<CompositionalMultiphaseFVM>} solver. A typo in either copy silently drops the region from the solver's target set; \texttt{xmllint} does not catch this.
\item \textbf{Solver--target coherence with \texttt{Tasks} and \texttt{Outputs}.} The \texttt{<PackCollection>} in \texttt{<Tasks>} dereferences \texttt{ElementRegions/region/cellBlock}; the \texttt{<TimeHistory>} output references \texttt{/Tasks/phaseVolumeFractionCollection}; the \texttt{<PeriodicEvent>} in the benchmark file targets \texttt{/Tasks/phaseVolumeFractionCollection} and \texttt{/Outputs/timeHistoryOutput}. Three sections (\texttt{Tasks}, \texttt{Outputs}, \texttt{Events}) must agree on the same name space.
\item \textbf{Constitutive references must point to declared blocks.} \texttt{materialList="{ fluid, rock, relperm }"} on the \texttt{<CellElementRegion>} requires that each of \texttt{fluid}, \texttt{rock}, \texttt{relperm} is the \texttt{name} of a present \texttt{<Constitutive>} child; \texttt{rock} in turn references \texttt{nullSolid}, \texttt{rockPorosity}, \texttt{rockPerm} via its constructor parameters, which must also exist.
\item \textbf{Geometry sets feed back into BCs.} The \texttt{<Box>} elements in \texttt{<Geometry>} declare the named sets \texttt{source} and \texttt{sink}, which are then referenced by \texttt{setNames} in the \texttt{<FieldSpecifications>} and the \texttt{<SourceFlux>}. A misspelled set name does not throw a parse error; it produces a deck that runs but with a missing source term.
\item \textbf{External tabulated functions.} \texttt{tableFiles="{ buckleyLeverett\_table/pvdg.txt, buckleyLeverett\_table/pvtw.txt }"} on \texttt{<DeadOilFluid>} requires those files to exist on disk relative to the deck location. The test harness copies the \texttt{buckleyLeverett\_table/} directory alongside the deck; without it, the deck parses but fails to construct the fluid model at runtime.
\end{itemize}

These constraints are what make GEOS XML behave as a small DSL rather than as a structured form: an authoring agent must keep five name spaces (regions, materials, sets, tasks, outputs) mutually consistent across blocks while also producing schema-valid attribute values. The \texttt{missing\_block} category our adapters reliably fix (\S\ref{subsec:bottleneck-results}) corresponds to whole-block omissions; the \texttt{bad\_attribute\_value} category they do not fix corresponds to typos and hallucinations in the cross-section name spaces above.

\section{Efficiency table}
\label{app:results}

\begin{table}[h]
  \caption{Mean tool calls and wall-clock per task ($n=3$ seeds). Wall is per-task seconds. \texttt{deepseek-v4-flash} output tokens are not exposed by the API; input tokens are dominated by cache-read of the system prompt.}
  \label{tab:efficiency}
  \centering
  \small
  \begin{tabular}{lcccc}
    \toprule
    \textbf{Cell} & \textbf{tools/task (val)} & \textbf{wall s (val)} & \textbf{tools/task (held-out-eval)} & \textbf{wall s (held-out-eval)} \\
    \midrule
    Vanilla  & 81.5 & 359 & 90.5 & 417 \\
    X+M  & 79.6 & 337 & 75.0 & 340 \\
    S+X  & 83.3 & 348 & 74.7 & 345 \\
    S+X+M  & 71.0 & 326 & 82.9 & 358 \\
    SE-prose & 62.7 & 326 & 70.9 & 362 \\
    SE  & 68.9 & 321 & 97.4 & 390 \\
    \bottomrule
  \end{tabular}
\end{table}

\section{Memory cheatsheet excerpt}
\label{app:cheatsheet}

The \textbf{M} factor delivers the file \texttt{plugin/memory\_primer\_m1u.md} via Claude Code's \texttt{--append-system-prompt}. The cheatsheet is a compact (775-token) enumerated reference: solver families paired with concrete XML element names, common constitutive-model class names, and a short list of anti-patterns. We reproduce a representative excerpt below; the full file is included in the supplementary material.

\begin{quote}\small
\textbf{Solver Selection \& Physics Routing}\\[2pt]
\noindent\resizebox{0.98\linewidth}{!}{%
\scriptsize
\setlength{\tabcolsep}{4pt}
\begin{tabular}{@{}p{0.18\linewidth} p{0.36\linewidth} p{0.42\linewidth}@{}}
\toprule
Physics family & Primary solver & Key supporting elements \\
\midrule
Hydrofracture           & \texttt{Hydrofracture}                  & \texttt{SurfaceGenerator}, \texttt{ParallelPlatesPermeability} \\
Solid mechanics         & \texttt{SolidMechanicsLagrangianFEM}    & \texttt{ElasticIsotropic}, \texttt{DruckerPrager} \\
Poromechanics           & \texttt{SinglePhasePoromechanics}       & requires \texttt{flowSolverName}, \texttt{solidSolverName} \\
Thermal flow            & \texttt{SinglePhaseThermalFVM}          & \texttt{SinglePhaseThermalConductivity}, \texttt{SolidInternalEnergy} \\
Multiphase flow         & \texttt{CompositionalMultiphaseFVM}     & \texttt{DeadOilFluid}, \texttt{CompositionalMultiphaseWell}, \texttt{SourceFlux} \\
Contact / interfaces    & \texttt{SolidMechanicsLagrangeContact}  & \texttt{SurfaceGenerator}, \texttt{EmbeddedSurface} \\
\bottomrule
\end{tabular}}
\\[6pt]
\textbf{Common anti-patterns (do NOT use).}\\
\begin{itemize}
\item Do NOT use \texttt{<FractureModel>} or \texttt{<HydraulicFractureSolver>}; these are hallucinated. Use \texttt{Hydrofracture} and \texttt{SurfaceGenerator}.
\item Do NOT use \texttt{<ContactSolver>}; use \texttt{SolidMechanicsLagrangeContact}.
\item Do NOT use \texttt{<FluidProperties>} as a wrapper; fluid models like \texttt{DeadOilFluid} or \texttt{CompressibleSinglePhaseFluid} are defined directly within the constitutive section.
\item Do NOT invent generic attribute names like \texttt{toughness} on the solver; verify benchmark-specific attributes (e.g., \texttt{kgdToughnessDominated}) via retrieval.
\end{itemize}
\textbf{Actionable tips.}\\
\begin{itemize}
\item Coupling: when using \texttt{SinglePhasePoromechanics}, both a flow solver and a solid solver must be defined and referenced by name.
\item Boundary conditions: confirm whether a \texttt{SourceFlux} requires a scalar \texttt{scale} or a \texttt{TableFunction} for time-dependent injection.
\end{itemize}
\end{quote}

The cheatsheet is intentionally a vocabulary dump, not a policy: it does not tell the agent which physics family the current task requires, only which element names to use once that decision is made. Section~\ref{subsec:bottleneck-results} traces the M-factor lift to reduced \texttt{missing\_block} errors (the cheatsheet enumerates the canonical sections and the families' member elements) and the residual \texttt{bad\_attribute\_value} errors to the cross-section name-consistency problem the cheatsheet does not address.


\section{Representative trajectory}
\label{app:case-study}

To make the agent's working pattern concrete, we reproduce an abridged transcript from the \texttt{buckleyLeverettProblem} task under the \emph{X+M} cell (\texttt{deepseek-v4-flash}, the canonical headline backbone), matching the deck shown in App.~\ref{app:geos-example}. The trajectory shows (i) initial planning, (ii) on-disk example-library reads (the X+M cell does not have RAG; analogous behaviour with retrieval calls is described in \S\ref{subsec:efficiency}), (iii) XML emission, (iv) the agent's voluntary \texttt{xmllint} validation calls (the X factor; this cell has no stop hook), and (v) the final TreeSim breakdown by section.

\promptinput{assets/trajectory_buckleyleverett_xm.txt}

\paragraph{What this trajectory illustrates.} The agent's translation strategy is example-driven: the first move is to locate a structurally analogous deck (\texttt{compositionalMultiphaseFlow/dbc/buckleyLeverett\_1d/}) under the same physics family and use it as a structural template, then re-parameterise. This is the same strategy the human-baseline analysis (\S\ref{subsec:human-baseline}) finds the human authors approximating via Sphinx browsing under time pressure. The voluntary \texttt{xmllint} calls (the X factor) catch one schema violation mid-trajectory (a missing \texttt{component} attribute on a \texttt{FieldSpecification}), which the agent fixes before the next emission; this is the on-trajectory analogue of the missing-block-resilience the bottleneck analysis attributes to schema-aware adapters (\S\ref{subsec:bottleneck-results}). The residual TreeSim loss (0.917 versus the per-task X+M mean of $\sim$0.92 on this task family, see Table~\ref{tab:per-task-icl10} for the held-out-eval analogue) sits inside the \texttt{bad\_attribute\_value} / \texttt{partial\_implementation} regime: a paired \texttt{sinkTermComposition\_water} \texttt{<FieldSpecification>} is elided alongside its gas counterpart, and one \texttt{PeriodicEvent} for the time-history collection is dropped from the events block. Both are the kind of completeness errors a stop-hook (S) running coverage checks would have caught; this cell does not have S enabled, which is consistent with the \S\ref{subsec:openfoam-transfer} finding that S is the dominant transferable-reliability factor.


%%%%%%%%%%%%%%%%%%%%%%%%%%%%%%%%%%%%%%%%%%%%%%%%%%%%%%%%%%%%

\newpage
\input{checklist.tex}


\end{document}